%% 3-17 偶然のばらつき：確率と正規分布（第12回：組 V3-8 が書いた）
%% 数（確率・組合せの数・サイコロの和の分布・±σ の割合）は、すべて Python で計算して確かめた（v3-8w/d17.py）。
%% Σ・n!・e は使わない（5-10・5-17・5-14）。分散・σ は 3-14 と同じ「n でわる」約束。
\makeatletter
\newcommand\bkthrrline[1]{\par\addvspace{3pt}{\centering #1\par}\addvspace{3pt}}
\newcommand\bkthrrstep[1]{\par\addvspace{6pt}\noindent{\sffamily\bfseries\color{bkDeep}#1}\par\nobreak\vspace{2pt}}
% サイコロの目  \bkthrrdie{(x,y)}{目の数}{大きさmm}
\newcommand\bkthrrdie[3]{\begin{scope}[shift={#1}, scale=#3/10]
  \filldraw[fill=white, draw=bkInk, line width=0.6pt, rounded corners=1.6mm] (-5,-5) rectangle (5,5);
  \ifcase#2\or \fill[bkRed] (0,0) circle (1.5);
  \or \fill[bkInk] (-2.4,2.4) circle (1) (2.4,-2.4) circle (1);
  \or \fill[bkInk] (-2.6,2.6) circle (1) (0,0) circle (1) (2.6,-2.6) circle (1);
  \or \fill[bkInk] (-2.4,2.4) circle (1) (2.4,2.4) circle (1) (-2.4,-2.4) circle (1) (2.4,-2.4) circle (1);
  \or \fill[bkInk] (-2.6,2.6) circle (1) (2.6,2.6) circle (1) (0,0) circle (1) (-2.6,-2.6) circle (1) (2.6,-2.6) circle (1);
  \or \fill[bkInk] (-2.4,2.9) circle (1) (2.4,2.9) circle (1) (-2.4,0) circle (1) (2.4,0) circle (1) (-2.4,-2.9) circle (1) (2.4,-2.9) circle (1);
  \fi
\end{scope}}
% 確率のぼうグラフ（小）  \bkthrrbars{題}{(x,y) の列}{xmin}{xmax}{ymax}
\newcommand\bkthrrbars[5]{%
  \begin{tikzpicture}
  \begin{axis}[bkplot, axis lines=left, width=36mm, height=32mm, xmin=#3, xmax=#4, ymin=0, ymax=#5, ytick=\empty, xtick=\empty, grid=none,
    title={\sffamily\footnotesize #1}, title style={yshift=-1mm}]
    \addplot[ybar, bar width=0.62, fill=bkBlue!55, draw=bkBlue!80!black, line width=0.3pt] coordinates {#2};
  \end{axis}
  \end{tikzpicture}}
\makeatother

\bkchapter{3}{17}{\R{偶然}{ぐう|ぜん}のばらつき：\\ \R{確率}{かく|りつ}と\R{正規分布}{せい|き|ぶん|ぷ}}{\bklevel{3}}{\R{偶然}{ぐう|ぜん}のばらつき：\R{確率}{かく|りつ}と\R{正規分布}{せい|き|ぶん|ぷ}}
\bkmichishirube{2-5,3-13,3-14}{5-17,9-4}

\begin{goal}
\begin{itemize}[label={\color{\bkTT}●}, leftmargin=1.6\zw, itemsep=0pt, topsep=1pt]
\item \bklevel{2}\enspace \R{表}{ひょう}や\R{樹形図}{じゅ|けい|ず}で\R{場合}{ば|あい}を\R{数}{かぞ}えて、\R{確率}{かく|りつ} $P(A)$ と\R{期待値}{き|たい|ち}を\R{求}{もと}められる。
\item \bklevel{3}\enspace コイン $n$ \R{枚}{まい}で\R{表}{おもて}が $k$ \R{枚}{まい}になる\R{場合}{ば|あい}の\R{数}{かず}を、パスカルの\R{三角形}{さん|かく|けい}で\R{求}{もと}められる。
\item \bklevel{3}\enspace たくさん\R{足}{た}すとつりがね\R{形}{かたち}（\R{正規分布}{せい|き|ぶん|ぷ}）になることを\R{知}{し}り、\R{平均}{へい|きん}±$\sigma$ に\R{約}{やく}68\%、±$2\sigma$ に\R{約}{やく}95\% を\R{使}{つか}える。
\end{itemize}
\end{goal}

\begin{hatena}
\noindent\begin{minipage}[c]{0.5\linewidth}
サイコロ1\R{個}{こ}では、1から6のどの\R{目}{め}も\R{同}{おな}じくらい\R{出}{で}ます（グラフは\R{平}{たい}ら）。\par
では、2\R{個}{こ}、3\R{個}{こ}、10\R{個}{こ}とふって、\R{目}{め}を\R{足}{た}すと？\quad \R{足}{た}した\R{数}{かず}をたくさん\R{記録}{き|ろく}すると、\R{右}{みぎ}のように、まん\R{中}{なか}が\R{高}{たか}い「つりがね」の\R{形}{かたち}になっていきます。\par
\R{平}{たい}らなものを\R{足}{た}すと、なぜ\R{山}{やま}になるのでしょう。
\end{minipage}\hfill
\begin{minipage}[c]{0.47\linewidth}\centering
\begin{tikzpicture}[x=1mm, y=1mm]
  \foreach \k in {1,...,6} {\fill[bkBlue!55] ({(\k-1)*3.4},0) rectangle ({(\k-1)*3.4+2.6},8);}
  \node[font=\sffamily\footnotesize, anchor=south] at (10,13) {1\R{個}{こ}};
  \foreach \s/\c in {2/1,3/2,4/3,5/4,6/5,7/6,8/5,9/4,10/3,11/2,12/1} {\fill[bkBlue!55] ({26+(\s-2)*2},0) rectangle ({26+(\s-2)*2+1.6},{\c*1.6});}
  \node[font=\sffamily\footnotesize, anchor=south] at (37,13) {2\R{個}{こ}の\R{和}{わ}};
  \foreach \s/\c in {3/1,4/3,5/6,6/10,7/15,8/21,9/25,10/27,11/27,12/25,13/21,14/15,15/10,16/6,17/3,18/1} {\fill[bkBlue!55] ({(\s-3)*1.4},-24) rectangle ({(\s-3)*1.4+1.1},{-24+\c*0.36});}
  \node[font=\sffamily\footnotesize, anchor=south] at (11,-11) {3\R{個}{こ}の\R{和}{わ}};
  \foreach \s/\c in {20/14,21/24,22/40,23/63,24/95,25/137,26/190,27/254,28/326,29/405,30/485,31/561,32/629,33/682,34/715,35/727,36/715,37/682,38/629,39/561,40/485,41/405,42/326,43/254,44/190,45/137,46/95,47/63,48/40,49/24,50/14} {\fill[bkBlue!65] ({26+(\s-20)*0.72},-24) rectangle ({26+(\s-20)*0.72+0.6},{-24+\c*0.0135});}
  \node[font=\sffamily\footnotesize, anchor=south] at (37,-11) {10\R{個}{こ}の\R{和}{わ}};
  \draw[bkInk, line width=0.5pt] (-1,0) -- (21,0) (25,0) -- (49,0) (-1,-24) -- (23,-24) (25,-24) -- (49,-24);
\end{tikzpicture}
\end{minipage}
\end{hatena}

2-6・3-13・3-14では、はかった\R{数}{かず}には「ずれ」があり、\R{何回}{なん|かい}もはかるとちらばることを\R{学}{まな}びました。ずれの\R{大}{おお}きさは、\R{標準偏差}{ひょう|じゅん|へん|さ} $\sigma$ で\R{表}{あらわ}せました。この\R{章}{しょう}では、そのずれの\R{正体}{しょう|たい}である「\R{偶然}{ぐう|ぜん}」を、\R{数}{かず}で\R{考}{かんが}えます。

\R{偶然}{ぐう|ぜん}に\R{起}{お}きることは、1\R{回}{かい}1\R{回}{かい}は\R{予想}{よ|そう}できません。コインを\R{投}{な}げて\R{表}{おもて}が\R{出}{で}るか\R{裏}{うら}が\R{出}{で}るかは、\R{投}{な}げるまでわかりません。ところが、\textbf{たくさんくり\R{返}{かえ}すと、きまりが\R{見}{み}えてきます}。そのきまりを\R{表}{あらわ}す\R{数}{かず}が、\textbf{\R{確率}{かく|りつ}}です。この\R{章}{しょう}の\R{最後}{さい|ご}には、3-14で\R{出}{で}てきた「\R{平均}{へい|きん}±$\sigma$ にだいたい7\R{割}{わり}」の「7\R{割}{わり}」が、どこから\R{来}{く}るのかもわかります。

なお、この\R{章}{しょう}では「ほぼ\R{等}{ひと}しい」を\R{表}{あらわ}す\R{記号}{き|ごう} $\approx$ を\R{使}{つか}います（くわしくは5-8）。たとえば $\dfrac{1}{3}\approx0.33$ です。

%% ---------------- やってみよう・思い出そう ----------------
\begin{yattemiyou}[コインを\R{何回}{なん|かい}も\R{投}{な}げて、\R{表}{おもて}の\R{割合}{わり|あい}を\R{調}{しら}べよう]
\begin{enumerate}[label={\bkcircled[bkGreen]{\arabic*}}, leftmargin=2.4\zw, itemsep=1pt, topsep=2pt]
\item 10\R{円玉}{えん|だま}を1\R{枚}{まい}\R{投}{な}げて、\R{表}{おもて}（\R{平等院}{びょう|どう|いん}の\R{絵}{え}がある\R{面}{めん}）か\R{裏}{うら}かを\R{記録}{き|ろく}しよう。
\item 10\R{回}{かい}、20\R{回}{かい}、50\R{回}{かい}、100\R{回}{かい}\R{投}{な}げたところで、「\R{表}{おもて}が\R{出}{で}た\R{回数}{かい|すう}÷\R{投}{な}げた\R{回数}{かい|すう}」を\R{計算}{けい|さん}しよう（\R{表}{おもて}の\R{割合}{わり|あい}、2-5）。
\item \R{表計算}{ひょう|けい|さん}なら、A\R{列}{れつ}に \texttt{=RANDBETWEEN(0,1)} と\R{入}{い}れて1000\R{行}{ぎょう}コピーすると、0か1が\R{出}{で}ます。1を\R{表}{おもて}として、1000\R{回分}{かい|ぶん}を\R{一度}{いち|ど}にためせます（2-18）。
\item \R{投}{な}げる\R{回数}{かい|すう}をふやすと、\R{表}{おもて}の\R{割合}{わり|あい}はどうなっていくかな？
\end{enumerate}
\end{yattemiyou}

\begin{omoidasou}
\bkrecall{2-5}{\R{割合}{わり|あい}＝くらべる\R{量}{りょう}÷もとにする\R{量}{りょう}。たとえば20\R{回}{かい}\R{中}{なか}12\R{回}{かい}なら $12\div20=0.6$}
\bkrecall{2-7}{\R{同}{おな}じ\R{数}{かず}をくり\R{返}{かえ}しかけるのが\R{累乗}{るい|じょう}。$2^3=2\times2\times2=8$、$2^{10}=1024$}
\bkrecall{3-13}{\R{度}{ど}\R{数}{かず}÷\R{合計}{ごう|けい}の\R{数}{かず}＝\R{相対度数}{そう|たい|ど|すう}。ヒストグラムで、データのちらばりを\R{見}{み}る。}
\bkrecall{3-14}{\R{標準偏差}{ひょう|じゅん|へん|さ} $\sigma$＝\R{偏差}{へん|さ}の2\R{乗}{じょう}の\R{平均}{へい|きん}（\R{分散}{ぶん|さん}、$n$ でわる）の $\sqrt{\phantom{x}}$。\R{平均}{へい|きん}±$\sigma$ に、データのだいたい7\R{割}{わり}が\R{入}{はい}る。}
\end{omoidasou}

\bkfukidashi{mina}{10\R{回}{かい}\R{投}{な}げたら\R{表}{おもて}が7\R{回}{かい}も\R{出}{で}たよ。コインがかたよっているのかな？}

%% ---------------- 数学の道具 ----------------
\bksection{doguu}{\R{数学}{すう|がく}の\R{道具}{どう|ぐ}}

\bksubsection{1}{たくさんくり\R{返}{かえ}すと、\R{割合}{わり|あい}が\R{落}{お}ち\R{着}{つ}く}

\leavevmode\label{bk:3-17:tokei}ソラさんとミナさんが、\R{表計算}{ひょう|けい|さん}でコインを1000\R{回}{かい}ずつ\R{投}{な}げました。\R{投}{な}げた\R{回数}{かい|すう}ごとに、それまでの\R{表}{おもて}の\R{割合}{わり|あい}をグラフにすると、\R{図}{ず}\ref{bk:3-17:fig-coin}のようになりました。

\begin{bkfigure}
\begin{tikzpicture}
\begin{axis}[bkplot, axis lines=left, width=124mm, height=66mm, xmin=0, xmax=1000, ymin=0, ymax=1.05, xtick={0,100,...,1000}, ytick={0,0.25,0.5,0.75,1}, xlabel={\R{投}{な}げた\R{回数}{かい|すう}}, ylabel={\R{表}{おもて}の\R{割合}{わり|あい}}, yticklabel style={/pgf/number format/fixed, /pgf/number format/precision=2}]
  \addplot[bkGray, dashed, line width=0.8pt, domain=0:1000] {0.5};
  \addplot[bkBlue, line width=1.1pt] coordinates {(1,1.000) (2,0.500) (3,0.333) (4,0.500) (5,0.600) (6,0.667) (7,0.714) (8,0.625) (9,0.556) (10,0.500) (11,0.455) (12,0.417) (13,0.462) (14,0.500) (15,0.533) (16,0.500) (17,0.529) (18,0.500) (19,0.526) (20,0.500) (25,0.560) (30,0.500) (35,0.543) (40,0.550) (45,0.578) (50,0.540) (55,0.545) (60,0.533) (65,0.538) (70,0.543) (75,0.560) (80,0.575) (85,0.576) (90,0.567) (95,0.579) (100,0.570) (125,0.544) (150,0.507) (175,0.503) (200,0.520) (225,0.520) (250,0.520) (275,0.527) (300,0.510) (325,0.517) (350,0.509) (375,0.512) (400,0.520) (425,0.513) (450,0.518) (475,0.514) (500,0.506) (525,0.501) (550,0.504) (575,0.499) (600,0.500) (625,0.502) (650,0.506) (675,0.505) (700,0.503) (725,0.513) (750,0.512) (775,0.514) (800,0.510) (825,0.508) (850,0.508) (875,0.509) (900,0.509) (925,0.507) (950,0.504) (975,0.503) (1000,0.501)};
  \addplot[bkOrange, line width=1.1pt] coordinates {(1,0.000) (2,0.000) (3,0.000) (4,0.000) (5,0.000) (6,0.000) (7,0.143) (8,0.250) (9,0.222) (10,0.200) (11,0.182) (12,0.250) (13,0.308) (14,0.357) (15,0.333) (16,0.312) (17,0.353) (18,0.389) (19,0.421) (20,0.400) (25,0.400) (30,0.467) (35,0.429) (40,0.400) (45,0.444) (50,0.460) (55,0.491) (60,0.500) (65,0.492) (70,0.500) (75,0.493) (80,0.487) (85,0.482) (90,0.478) (95,0.495) (100,0.500) (125,0.504) (150,0.500) (175,0.480) (200,0.455) (225,0.489) (250,0.480) (275,0.491) (300,0.493) (325,0.489) (350,0.489) (375,0.493) (400,0.492) (425,0.496) (450,0.498) (475,0.505) (500,0.514) (525,0.522) (550,0.522) (575,0.525) (600,0.528) (625,0.525) (650,0.528) (675,0.530) (700,0.523) (725,0.516) (750,0.515) (775,0.514) (800,0.512) (825,0.513) (850,0.518) (875,0.517) (900,0.518) (925,0.519) (950,0.519) (975,0.520) (1000,0.521)};
  \node[font=\sffamily\footnotesize, text=bkBlue!80!black, anchor=west] at (axis cs:380,0.92) {ソラさん（1000\R{回}{かい}で0.501）};
  \node[font=\sffamily\footnotesize, text=bkOrange!80!black, anchor=west] at (axis cs:380,0.82) {ミナさん（1000\R{回}{かい}で0.521）};
\end{axis}
\end{tikzpicture}
\bkfigcaption{\R{表}{おもて}の\R{割合}{わり|あい}の\R{変}{か}わり\R{方}{かた}。はじめは\R{大}{おお}きくゆれるが、\R{回数}{かい|すう}がふえると0.5の\R{近}{ちか}くに\R{落}{お}ち\R{着}{つ}く。}\label{bk:3-17:fig-coin}
\end{bkfigure}

はじめの10\R{回}{かい}くらいまでは、\R{割合}{わり|あい}は\R{大}{おお}きくゆれます。ミナさんは、はじめに6\R{回}{かい}\R{続}{つづ}けて\R{裏}{うら}が\R{出}{で}たので、\R{割合}{わり|あい}は0でした。でも、\R{何百回}{なん|びゃっ|かい}と\R{投}{な}げていくと、2\R{人}{にん}とも\R{割合}{わり|あい}が0.5の\R{近}{ちか}くに\R{落}{お}ち\R{着}{つ}いていきます。1\R{回}{かい}1\R{回}{かい}はでたらめなのに、\textbf{たくさんくり\R{返}{かえ}すと、\R{割合}{わり|あい}はある\R{数}{かず}に\R{近}{ちか}づく}のです。

このように、たくさんくり\R{返}{かえ}したときに\R{割合}{わり|あい}が\R{近}{ちか}づいていく\R{数}{かず}を、そのことがらの\R{起}{お}こる\textbf{\R{確率}{かく|りつ}}といいます。\bkidx{かくりつ}{\R{確率}{かく|りつ}}コインの\R{表}{おもて}が\R{出}{で}る\R{確率}{かく|りつ}は0.5、つまり $\dfrac{1}{2}$ です。\R{実験}{じっ|けん}で、くり\R{返}{かえ}した\R{回数}{かい|すう}から\R{求}{もと}める\R{確率}{かく|りつ}なので、\textbf{\R{統計的}{とう|けい|てき}な\R{確率}{かく|りつ}}ともいいます。

\textbf{\R{画}{が}びょうの\R{確率}{かく|りつ}}\quad \R{画}{が}びょうを\R{投}{な}げると、\R{針}{はり}が\R{上}{うえ}を\R{向}{む}く「\R{上}{うえ}\R{向}{む}き」と、\R{横}{よこ}にたおれる「\R{横}{よこ}\R{向}{む}き」があります。\R{形}{かたち}が\R{左右}{さ|ゆう}で\R{同}{おな}じではないので、どちらが\R{出}{で}やすいかは、\R{考}{かんが}えるだけではわかりません。\R{実験}{じっ|けん}するしかありません。

\begin{bkcenter}
\begin{bktable}{colspec={Q[l,m,wd=30mm] *{5}{Q[c,m,wd=14mm]}}, rowsep=1pt, cells={font=\small}}
\R{投}{な}げた\R{回数}{かい|すう} & 100 & 200 & 500 & 1000 & 2000 \\
\R{上}{うえ}\R{向}{む}きの\R{回数}{かい|すう} & 66 & 127 & 309 & 621 & 1238 \\
\R{上}{うえ}\R{向}{む}きの\R{割合}{わり|あい} & 0.660 & 0.635 & 0.618 & 0.621 & 0.619 \\
\end{bktable}
\bktabcaption{\R{画}{が}びょうを\R{投}{な}げた\R{結果}{けっ|か}（ある1つの\R{画}{が}びょう）}
\end{bkcenter}
\R{割合}{わり|あい}は0.62くらいに\R{落}{お}ち\R{着}{つ}いています。この\R{画}{が}びょうが\R{上}{うえ}\R{向}{む}きになる\R{確率}{かく|りつ}は、およそ0.62と\R{考}{かんが}えられます。2000\R{回}{かい}\R{投}{な}げたあとで、もう1000\R{回}{かい}\R{投}{な}げるなら、\R{上}{うえ}\R{向}{む}きはおよそ $1000\times0.62=620$ \R{回}{かい}と\R{予想}{よ|そう}できます。

\bkhatenaref に\R{向}{む}けて：ミナさんの「10\R{回}{かい}で\R{表}{おもて}が7\R{回}{かい}」も、ふしぎなことではありません。10\R{回}{かい}くらいでは、\R{割合}{わり|あい}は\R{大}{おお}きくゆれるからです。どのくらいゆれるのかは、③・⑤で\R{計算}{けい|さん}します。

\begin{column}{\R{天気予報}{てん|き|よ|ほう}の「\R{降水確率}{こう|すい|かく|りつ}30\%」}
\R{天気予報}{てん|き|よ|ほう}の「あしたの\R{降水確率}{こう|すい|かく|りつ}は30\%」は、どんな\R{意味}{い|み}でしょう。「あしたは、\R{1日}{いちにち}のうち30\%の\R{時間}{じ|かん}だけ\R{雨}{あめ}がふる」でも、「30\%の\R{地域}{ち|いき}で\R{雨}{あめ}がふる」でもありません。\par
「\R{今回}{こん|かい}と\R{同}{おな}じような\R{天気図}{てん|き|ず}の\R{日}{ひ}が、\R{過去}{か|こ}に100\R{回}{かい}あったとすると、そのうち\R{約}{やく}30\R{回}{かい}は、1\,mm\R{以上}{い|じょう}の\R{雨}{あめ}がふった」という\R{意味}{い|み}です。①の\R{統計的}{とう|けい|てき}な\R{確率}{かく|りつ}と\R{同}{おな}じ\R{考}{かんが}え\R{方}{かた}で、たくさんの\R{過去}{か|こ}のデータの\R{割合}{わり|あい}から\R{決}{き}めています。30\%なら、10\R{回}{かい}のうち3\R{回}{かい}くらいは\R{雨}{あめ}。かさを\R{持}{も}っていくかどうかは、あなたが\R{決}{き}めましょう。
\end{column}

\bksubsection{2}{\R{場合}{ば|あい}を\R{数}{かぞ}えて、\R{確率}{かく|りつ}を\R{求}{もと}める}

\leavevmode\label{bk:3-17:kazoeru}コインやサイコロは、\R{実験}{じっ|けん}しなくても\R{確率}{かく|りつ}がわかります。

\bkthrrstep{\R{同}{おな}じように\R{確}{たし}からしい}
\noindent\begin{minipage}[c]{0.66\linewidth}
きちんと\R{作}{つく}られたサイコロは、どの\R{面}{めん}も\R{同}{おな}じ\R{形}{かたち}・\R{同}{おな}じ\R{重}{おも}さです。だから、1から6のどの\R{目}{め}も、\R{出}{で}やすさは\R{同}{おな}じだと\R{考}{かんが}えられます。このように、どれが\R{起}{お}こることも\R{同}{おな}じ\R{程度}{てい|ど}に\R{期待}{き|たい}できるとき、「\textbf{\R{同様}{どう|よう}に\R{確}{たし}からしい}」といいます。
\end{minipage}\hfill
\begin{minipage}[c]{0.3\linewidth}\centering
\begin{tikzpicture}[x=1mm, y=1mm]
  \foreach \k in {1,...,6} {{\bkthrrdie{({(\k-1)*7},0)}{\k}{6}}}
\end{tikzpicture}
\end{minipage}

\bkthrrstep{\R{確率}{かく|りつ}を\R{数}{かず}で\R{求}{もと}める}
\R{同様}{どう|よう}に\R{確}{たし}からしい\R{場合}{ば|あい}が\R{全部}{ぜん|ぶ}で $n$ \R{通}{とお}りあり、そのうち、ことがら $A$ が\R{起}{お}こる\R{場合}{ば|あい}が $a$ \R{通}{とお}りあるとき、$A$ の\R{起}{お}こる\R{確率}{かく|りつ}は
\bkthrrline{$P(A)=\dfrac{a}{n}$\qquad（$A$ の\R{起}{お}こる\R{場合}{ば|あい}の\R{数}{かず}）÷（\R{全部}{ぜん|ぶ}の\R{場合}{ば|あい}の\R{数}{かず}）}
です。$P$ は\R{英語}{えい|ご}の probability（\R{確率}{かく|りつ}）の\R{頭文字}{かしら|も|じ}で、$P(A)$ は「ピー・エー」と\R{読}{よ}みます。$f(x)$ と\R{同}{おな}じかっこの\R{書}{か}き\R{方}{かた}（3-3）で、「ことがら $A$ を\R{入}{い}れると、その\R{確率}{かく|りつ}が\R{出}{で}てくる」と\R{読}{よ}めます。

\noindent\begin{minipage}[t]{0.46\linewidth}\vspace{0pt}
たとえば、サイコロ1\R{個}{こ}で\R{偶数}{ぐう|すう}の\R{目}{め}が\R{出}{で}るのは、2、4、6の3\R{通}{とお}りなので、
\bkthrrline{（\R{偶数}{ぐう|すう}の\R{目}{め}が\R{出}{で}る\R{確率}{かく|りつ}）$=\dfrac{3}{6}=\dfrac{1}{2}$}
です。\R{必}{かなら}ず\R{起}{お}こること（6\R{以下}{い|か}の\R{目}{め}が\R{出}{で}る）は $\dfrac{6}{6}=1$、\R{決}{き}して\R{起}{お}こらないこと（7の\R{目}{め}が\R{出}{で}る）は $\dfrac{0}{6}=0$ です。\R{確率}{かく|りつ}は、いつも0\R{以上}{い|じょう}1\R{以下}{い|か}の\R{数}{かず}です。
\end{minipage}\hfill
\begin{minipage}[t]{67mm}\vspace{0pt}\bkkigocard{$P(A)$}{ピー・エー}{ことがら $A$ の\R{起}{お}こる\R{確率}{かく|りつ}。0から1までの\R{数}{かず}。}{probability}{3-17}\end{minipage}
\par\vspace{4pt}

\bkthrrstep{2\R{個}{こ}のサイコロの\R{和}{わ}は、\R{表}{ひょう}で\R{数}{かぞ}える}
\R{大小}{だい|しょう}2\R{個}{こ}のサイコロをふって、\R{目}{め}の\R{和}{わ}を\R{考}{かんが}えます。\R{大}{おお}が1で\R{小}{しょう}が2の\R{場合}{ば|あい}と、\R{大}{おお}が2で\R{小}{しょう}が1の\R{場合}{ば|あい}は、\R{別}{べつ}の\R{場合}{ば|あい}です。\R{全部}{ぜん|ぶ}の\R{場合}{ば|あい}を、6×6の\R{表}{ひょう}にします（1-13の\R{九九}{く|く}の\R{表}{ひょう}と\R{同}{おな}じ\R{形}{かたち}）。

\begin{bkcenter}
\begin{bktable}{colspec={Q[c,m,wd=12mm] *{6}{Q[c,m,wd=9mm]}}, rowsep=0.5pt, colsep=3pt, cells={font=\small}, column{1}={bg=bkTheme!12, font=\sffamily\bfseries\small}}
\R{大}{おお}＼\R{小}{しょう} & 1 & 2 & 3 & 4 & 5 & 6 \\
1 & 2 & 3 & 4 & 5 & 6 & \SetCell{bg=bkRed!15}7 \\
2 & 3 & 4 & 5 & 6 & \SetCell{bg=bkRed!15}7 & 8 \\
3 & 4 & 5 & 6 & \SetCell{bg=bkRed!15}7 & 8 & 9 \\
4 & 5 & 6 & \SetCell{bg=bkRed!15}7 & 8 & 9 & 10 \\
5 & 6 & \SetCell{bg=bkRed!15}7 & 8 & 9 & 10 & 11 \\
6 & \SetCell{bg=bkRed!15}7 & 8 & 9 & 10 & 11 & 12 \\
\end{bktable}
\bktabcaption{2\R{個}{こ}のサイコロの\R{目}{め}の\R{和}{わ}。\R{全部}{ぜん|ぶ}で $6\times6=36$ \R{通}{とお}り。\R{和}{わ}が7（\R{赤}{あか}）は6\R{通}{とお}り。}
\end{bkcenter}

\R{和}{わ}が7になるのは、ななめに\R{並}{なら}んだ6\R{通}{とお}りです。だから、\R{和}{わ}が7になる\R{確率}{かく|りつ}は $\dfrac{6}{36}=\dfrac{1}{6}$ です。\R{和}{わ}が2になるのは「\R{大}{おお}1・\R{小}{しょう}1」の1\R{通}{とお}りだけなので、$\dfrac{1}{36}$ です。\R{和}{わ}ごとの\R{場合}{ば|あい}の\R{数}{かず}を\R{数}{かぞ}えると、\R{次}{つぎ}のようになります。

\noindent\begin{minipage}[c]{0.5\linewidth}\centering
\begin{bktable}{colspec={Q[c,m,wd=12mm] *{6}{Q[c,m,wd=5.5mm]}}, rowsep=0.5pt, colsep=1.5pt, cells={font=\small}}
\R{和}{わ} & 2 & 3 & 4 & 5 & 6 & 7 \\
\R{場合}{ば|あい} & 1 & 2 & 3 & 4 & 5 & 6 \\
\R{和}{わ} & 8 & 9 & 10 & 11 & 12 & \\
\R{場合}{ば|あい} & 5 & 4 & 3 & 2 & 1 & \\
\end{bktable}
\bktabcaption{\R{和}{わ}ごとの\R{場合}{ば|あい}の\R{数}{かず}（\R{合計}{ごう|けい}36）}
\end{minipage}\hfill
\begin{minipage}[c]{0.46\linewidth}\centering
\begin{tikzpicture}
\begin{axis}[bkplot, axis lines=left, width=58mm, height=40mm, xmin=1.3, xmax=12.7, ymin=0, ymax=7, xtick={2,4,...,12}, ytick={0,2,4,6}, xlabel={\R{和}{わ}}, ylabel={\R{場合}{ば|あい}の\R{数}{かず}}, grid=none]
  \addplot[ybar, bar width=3.2mm, fill=bkBlue!55, draw=bkBlue!80!black] coordinates {(2,1) (3,2) (4,3) (5,4) (6,5) (7,6) (8,5) (9,4) (10,3) (11,2) (12,1)};
\end{axis}
\end{tikzpicture}
\end{minipage}

1\R{個}{こ}のサイコロでは、どの\R{目}{め}も\R{同}{おな}じ\R{確率}{かく|りつ}でした（\R{平}{たい}ら）。2\R{個}{こ}の\R{和}{わ}になると、7がいちばん\R{出}{で}やすい\R{山}{やま}の\R{形}{かたち}になります。まん\R{中}{なか}の7は、\R{作}{つく}り\R{方}{かた}（\R{大}{おお}と\R{小}{しょう}の\R{組}{くみ}）がいちばん\R{多}{おお}いからです。\bkhatenaref の\R{答}{こた}えの、1つめのかぎです。

\bksubsection{3}{コインの\R{表}{おもて}の\R{枚数}{まい|すう}を\R{数}{かぞ}える：パスカルの\R{三角形}{さん|かく|けい}}

\leavevmode\label{bk:3-17:pascal}コインを\R{何枚}{なん|まい}か\R{投}{な}げて、\R{表}{おもて}が\R{何枚}{なん|まい}\R{出}{で}るかを\R{考}{かんが}えます。

\bkthrrstep{\R{段}{だん}1　\R{樹形図}{じゅ|けい|ず}で、\R{全部}{ぜん|ぶ}の\R{場合}{ば|あい}を\R{書}{か}き\R{出}{だ}す}
コインを3\R{枚}{まい}、1\R{枚}{まい}めから\R{順}{じゅん}に\R{投}{な}げます。1\R{枚}{まい}ごとに、\R{表}{おもて}（お）か\R{裏}{うら}（う）の2\R{通}{とお}りに\R{分}{わ}かれていく\R{図}{ず}をかきます。\R{木}{き}の\R{枝}{えだ}のように\R{分}{わ}かれるので、\textbf{\R{樹形図}{じゅ|けい|ず}}といいます。

\begin{bkfigure}
\begin{tikzpicture}[x=1mm, y=1mm, every node/.style={font=\sffamily\small, inner sep=1.2pt}]
  \node (s) at (0,0) {};
  \fill[bkInk] (0,0) circle (0.7);
  \foreach \a/\ya in {お/16, う/-16} {
    \node[draw=bkInk, circle, minimum size=5mm] (A\a) at (20,\ya) {\a};
    \draw[bkInk, line width=0.6pt] (0.8,0) -- (A\a);
    \foreach \b/\yb in {お/8, う/-8} {
      \node[draw=bkInk, circle, minimum size=5mm] (B\a\b) at (40,{\ya+\yb}) {\b};
      \draw[bkInk, line width=0.6pt] (A\a) -- (B\a\b);
      \foreach \c/\yc in {お/4, う/-4} {
        \node[draw=bkInk, circle, minimum size=5mm] (C\a\b\c) at (60,{\ya+\yb+\yc}) {\c};
        \draw[bkInk, line width=0.6pt] (B\a\b) -- (C\a\b\c);
        \node[anchor=west] at (68,{\ya+\yb+\yc}) {\a\b\c};}}}
  \node[font=\sffamily\footnotesize, anchor=south] at (20,32) {1\R{枚}{まい}め};
  \node[font=\sffamily\footnotesize, anchor=south] at (40,32) {2\R{枚}{まい}め};
  \node[font=\sffamily\footnotesize, anchor=south] at (60,32) {3\R{枚}{まい}め};
  \foreach \y/\k in {28/3, 20/2, 12/2, 4/1, -4/2, -12/1, -20/1, -28/0} {
    \node[font=\sffamily\small, text=bkRed!80!black, anchor=west] at (82,\y-0.0) {\R{表}{おもて}\k \R{枚}{まい}};}
\end{tikzpicture}
\bkfigcaption{コイン3\R{枚}{まい}の\R{樹形図}{じゅ|けい|ず}。\R{全部}{ぜん|ぶ}で $2\times2\times2=8$ \R{通}{とお}り（お＝\R{表}{おもて}、う＝\R{裏}{うら}）。}\label{bk:3-17:fig-jukei}
\end{bkfigure}

\R{全部}{ぜん|ぶ}で $2\times2\times2=2^3=8$ \R{通}{とお}りです（2-7の\R{累乗}{るい|じょう}）。\R{表}{おもて}の\R{枚数}{まい|すう}ごとに\R{数}{かぞ}えると、

\begin{bkcenter}
\begin{bktable}{colspec={Q[l,m,wd=28mm] *{4}{Q[c,m,wd=14mm]}}, rowsep=0.5pt, cells={font=\small}}
\R{表}{おもて}の\R{枚数}{まい|すう} & 0\R{枚}{まい} & 1\R{枚}{まい} & 2\R{枚}{まい} & 3\R{枚}{まい} \\
\R{場合}{ば|あい}の\R{数}{かず} & 1 & 3 & 3 & 1 \\
\R{確率}{かく|りつ} & $\dfrac{1}{8}$ & $\dfrac{3}{8}$ & $\dfrac{3}{8}$ & $\dfrac{1}{8}$ \\
\end{bktable}
\bktabcaption{コイン3\R{枚}{まい}の\R{表}{おもて}の\R{枚数}{まい|すう}}
\end{bkcenter}
\R{表}{おもて}が2\R{枚}{まい}になるのは、「おおう」「おうお」「うおお」の3\R{通}{とお}りです。どの1\R{枚}{まい}が\R{裏}{うら}になるかで3\R{通}{とお}り、と\R{考}{かんが}えることもできます。1、3、3、1は、3-8の $(a+b)^3$ を\R{展開}{てん|かい}したときの\R{係数}{けい|すう}と\R{同}{おな}じ\R{数}{かず}です。

\bkthrrstep{\R{段}{だん}2　1\R{枚}{まい}ふやすと、どうなる？}
4\R{枚}{まい}にすると、\R{樹形図}{じゅ|けい|ず}は16\R{通}{とお}りになって、かくのが\R{大変}{たい|へん}です。そこで、\R{樹形図}{じゅ|けい|ず}をかかずに\R{数}{かぞ}える\R{方法}{ほう|ほう}を\R{考}{かんが}えます。

4\R{枚}{まい}で\R{表}{おもて}がちょうど2\R{枚}{まい}になるのは、4\R{枚}{まい}めのコインで2つに\R{分}{わ}けられます。
\begin{itemize}[leftmargin=1.4\zw, itemsep=1pt, topsep=2pt, label={\color{bkBlue}●}]
\item 4\R{枚}{まい}めが\R{裏}{うら}のとき：はじめの3\R{枚}{まい}で、\R{表}{おもて}が2\R{枚}{まい}。\R{前}{まえ}の\R{表}{ひょう}から3\R{通}{とお}り。
\item 4\R{枚}{まい}めが\R{表}{おもて}のとき：はじめの3\R{枚}{まい}で、\R{表}{おもて}が1\R{枚}{まい}。\R{前}{まえ}の\R{表}{ひょう}から3\R{通}{とお}り。
\end{itemize}
だから、$3+3=6$ \R{通}{とお}りです。\R{同}{おな}じように、4\R{枚}{まい}で\R{表}{おもて}が1\R{枚}{まい}は「3\R{枚}{まい}で1\R{枚}{まい}」＋「3\R{枚}{まい}で0\R{枚}{まい}」$=3+1=4$ \R{通}{とお}りです。

\textbf{きまり}：$n$ \R{枚}{まい}で\R{表}{おもて}が $k$ \R{枚}{まい}の\R{数}{かず}＝（$n-1$ \R{枚}{まい}で\R{表}{おもて}が $k$ \R{枚}{まい}の\R{数}{かず}）＋（$n-1$ \R{枚}{まい}で\R{表}{おもて}が $k-1$ \R{枚}{まい}の\R{数}{かず}）

\bkthrrstep{\R{段}{だん}3　\R{三角形}{さん|かく|けい}にならべる}
このきまりを\R{使}{つか}って、\R{上}{うえ}から1\R{段}{だん}ずつ\R{数}{かず}を\R{書}{か}いていきます。\R{各}{かく}\R{段}{だん}の\R{両}{りょう}はしは1です（\R{全部}{ぜん|ぶ}\R{裏}{うら}、\R{全部}{ぜん|ぶ}\R{表}{おもて}は1\R{通}{とお}りずつ）。そのほかの\R{数}{かず}は、\textbf{すぐ\R{上}{うえ}の\R{左右}{さ|ゆう}の2つの\R{数}{かず}をたした\R{数}{かず}}です。

\begin{bkfigure}
\begin{tikzpicture}[x=1mm, y=1mm]
  \node[font=\sffamily\footnotesize, text=bkGray, anchor=east] at (-6,0) {0\R{枚}{まい}};
  \node[font=\small] at (48.0,0) {1};
  \node[font=\sffamily\footnotesize, text=bkGray, anchor=east] at (-6,-8) {1\R{枚}{まい}};
  \node[font=\small] at (42.5,-8) {1};
  \node[font=\small] at (53.5,-8) {1};
  \node[font=\sffamily\footnotesize, text=bkGray, anchor=east] at (-6,-16) {2\R{枚}{まい}};
  \node[font=\small] at (37.0,-16) {1};
  \node[font=\small] at (48.0,-16) {2};
  \node[font=\small] at (59.0,-16) {1};
  \node[font=\sffamily\footnotesize, text=bkGray, anchor=east] at (-6,-24) {3\R{枚}{まい}};
  \node[font=\small] at (31.5,-24) {1};
  \node[font=\small] at (42.5,-24) {3};
  \node[font=\small] at (53.5,-24) {3};
  \node[font=\small] at (64.5,-24) {1};
  \node[font=\sffamily\footnotesize, text=bkGray, anchor=east] at (-6,-32) {4\R{枚}{まい}};
  \fill[bkYellow!40] (26.0,-32) circle (3.6);
  \node[font=\small] at (26.0,-32) {1};
  \fill[bkYellow!40] (37.0,-32) circle (3.6);
  \node[font=\small] at (37.0,-32) {4};
  \fill[bkYellow!40] (48.0,-32) circle (3.6);
  \node[font=\small] at (48.0,-32) {6};
  \fill[bkYellow!40] (59.0,-32) circle (3.6);
  \node[font=\small] at (59.0,-32) {4};
  \fill[bkYellow!40] (70.0,-32) circle (3.6);
  \node[font=\small] at (70.0,-32) {1};
  \node[font=\sffamily\footnotesize, text=bkGray, anchor=east] at (-6,-40) {5\R{枚}{まい}};
  \node[font=\small] at (20.5,-40) {1};
  \node[font=\small] at (31.5,-40) {5};
  \node[font=\small] at (42.5,-40) {10};
  \node[font=\small] at (53.5,-40) {10};
  \node[font=\small] at (64.5,-40) {5};
  \node[font=\small] at (75.5,-40) {1};
  \node[font=\sffamily\footnotesize, text=bkGray, anchor=east] at (-6,-48) {6\R{枚}{まい}};
  \node[font=\small] at (15.0,-48) {1};
  \node[font=\small] at (26.0,-48) {6};
  \node[font=\small] at (37.0,-48) {15};
  \node[font=\small] at (48.0,-48) {20};
  \node[font=\small] at (59.0,-48) {15};
  \node[font=\small] at (70.0,-48) {6};
  \node[font=\small] at (81.0,-48) {1};
  \node[font=\sffamily\footnotesize, text=bkGray, anchor=east] at (-6,-56) {7\R{枚}{まい}};
  \node[font=\small] at (9.5,-56) {1};
  \node[font=\small] at (20.5,-56) {7};
  \node[font=\small] at (31.5,-56) {21};
  \node[font=\small] at (42.5,-56) {35};
  \node[font=\small] at (53.5,-56) {35};
  \node[font=\small] at (64.5,-56) {21};
  \node[font=\small] at (75.5,-56) {7};
  \node[font=\small] at (86.5,-56) {1};
  \node[font=\sffamily\footnotesize, text=bkGray, anchor=east] at (-6,-64) {8\R{枚}{まい}};
  \node[font=\small] at (4.0,-64) {1};
  \node[font=\small] at (15.0,-64) {8};
  \node[font=\small] at (26.0,-64) {28};
  \node[font=\small] at (37.0,-64) {56};
  \node[font=\small] at (48.0,-64) {70};
  \node[font=\small] at (59.0,-64) {56};
  \node[font=\small] at (70.0,-64) {28};
  \node[font=\small] at (81.0,-64) {8};
  \node[font=\small] at (92.0,-64) {1};
  % たし算の矢印の例（5段めの10 = 4+6）
  \draw[bkRed, line width=0.7pt, -{Stealth[length=3pt]}] (38.2,-34.6) -- (41.4,-37.6);
  \draw[bkRed, line width=0.7pt, -{Stealth[length=3pt]}] (46.8,-34.6) -- (43.6,-37.6);
  \node[font=\sffamily\footnotesize, text=bkRed!80!black, anchor=west] at (84,-36) {\R{赤}{あか}い\R{矢印}{や|じるし}：$4+6=10$};
\end{tikzpicture}
\bkfigcaption{パスカルの\R{三角形}{さん|かく|けい}。$n$ \R{段}{だん}めの\R{左}{ひだり}から $k$ \R{番}{ばん}め（0から\R{数}{かぞ}える）が、コイン $n$ \R{枚}{まい}で\R{表}{おもて}が $k$ \R{枚}{まい}の\R{数}{かず}。4\R{枚}{まい}の\R{段}{だん}（\R{黄}{き}）は1、4、6、4、1。}\label{bk:3-17:fig-pascal}
\end{bkfigure}

この\R{三角形}{さん|かく|けい}を、\textbf{パスカルの\R{三角形}{さん|かく|けい}}といいます。\bkidx{ぱすかるのさんかくけい}{パスカルの\R{三角形}{さん|かく|けい}}370\R{年}{ねん}ほど\R{前}{まえ}に、フランスのパスカルがくわしく\R{研究}{けん|きゅう}したので、この\R{名前}{な|まえ}があります。

4\R{枚}{まい}の\R{段}{だん}は1、4、6、4、1です。3-14で、4\R{回}{かい}の\R{並}{なら}び\R{方}{かた}の\R{数}{かず}として\R{出}{で}てきた1、4、6、4、1と\R{同}{おな}じです。\R{表}{おもて}が0\R{枚}{まい}から4\R{枚}{まい}になる\R{数}{かず}を\R{全部}{ぜん|ぶ}たすと、$1+4+6+4+1=16=2^4$ で、\R{全部}{ぜん|ぶ}の\R{場合}{ば|あい}の\R{数}{かず}になります。どの\R{段}{だん}でも、たすと $2^n$ です（1\R{枚}{まい}ふやすごとに、\R{場合}{ば|あい}が2\R{倍}{ばい}に\R{分}{わ}かれるから）。

\bkthrrstep{\R{段}{だん}4　\R{組}{くみ}\R{合}{あ}せの\R{数}{かず}と、その\R{記号}{き|ごう}}
「4\R{枚}{まい}のうち、どの2\R{枚}{まい}が\R{表}{おもて}か」は、4\R{枚}{まい}から2\R{枚}{まい}を\textbf{\R{選}{えら}ぶ}\R{選}{えら}び\R{方}{かた}の\R{数}{かず}と\R{同}{おな}じです。\R{順番}{じゅん|ばん}を\R{考}{かんが}えずに、いくつかを\R{選}{えら}ぶ\R{選}{えら}び\R{方}{かた}を\textbf{\R{組}{くみ}\R{合}{あ}せ}といい、$n$ \R{個}{こ}から $r$ \R{個}{こ}を\R{選}{えら}ぶ\R{組}{くみ}\R{合}{あ}せの\R{数}{かず}を ${}_n\mathrm{C}_r$ と\R{書}{か}きます。$\mathrm{C}$ は\R{英語}{えい|ご}の combination（\R{組}{くみ}\R{合}{あ}せ）の\R{頭文字}{かしら|も|じ}です。パスカルの\R{三角形}{さん|かく|けい}の $n$ \R{段}{だん}めの\R{左}{ひだり}から $r$ \R{番}{ばん}め（0\R{番}{ばん}めから\R{数}{かぞ}える）が、${}_n\mathrm{C}_r$ です。

\noindent\begin{minipage}[t]{0.46\linewidth}\vspace{0pt}
${}_n\mathrm{C}_r$ は、かけ\R{算}{ざん}でも\R{計算}{けい|さん}できます。5\R{枚}{まい}から2\R{枚}{まい}を\R{選}{えら}ぶとき、1\R{枚}{まい}めの\R{選}{えら}び\R{方}{かた}は5\R{通}{とお}り、2\R{枚}{まい}めは\R{残}{のこ}りの4\R{通}{とお}りで、$5\times4=20$ \R{通}{とお}り。でも、「AとB」と「BとA」は\R{同}{おな}じ\R{組}{くみ}なので、どの\R{組}{くみ}も2\R{回}{かい}ずつ\R{数}{かぞ}えています（2\R{枚}{まい}の\R{並}{なら}べ\R{方}{かた} $2\times1=2$ \R{通}{とお}り）。
\end{minipage}\hfill
\begin{minipage}[t]{67mm}\vspace{0pt}\bkkigocard{${}_n\mathrm{C}_r$}{エヌ・シー・アール}{$n$ \R{個}{こ}から $r$ \R{個}{こ}を\R{選}{えら}ぶ\R{組}{くみ}\R{合}{あ}せの\R{数}{かず}。パスカルの\R{三角形}{さん|かく|けい}の\R{数}{かず}。}{combination}{3-17}\end{minipage}
\par\vspace{3pt}
だから、2でわります（\R{図}{ず}\ref{bk:3-17:fig-pairs}）。
\[ {}_5\mathrm{C}_2=\frac{5\times4}{2\times1}=\frac{20}{2}=10 \]
\begin{bkfigure}
\begin{tikzpicture}[x=1mm, y=1mm]
  \fill[bkBlue!18] (10.5,-2.6) rectangle (19.5,2.6); \node[font=\sffamily\footnotesize] at (15.0,0.0) {AB};
  \fill[bkBlue!18] (21.5,-2.6) rectangle (30.5,2.6); \node[font=\sffamily\footnotesize] at (26.0,0.0) {AC};
  \fill[bkBlue!18] (32.5,-2.6) rectangle (41.5,2.6); \node[font=\sffamily\footnotesize] at (37.0,0.0) {AD};
  \fill[bkBlue!18] (43.5,-2.6) rectangle (52.5,2.6); \node[font=\sffamily\footnotesize] at (48.0,0.0) {AE};
  \fill[bkOrange!18] (-0.5,-8.6) rectangle (8.5,-3.4); \node[font=\sffamily\footnotesize] at (4.0,-6.0) {BA};
  \fill[bkBlue!18] (21.5,-8.6) rectangle (30.5,-3.4); \node[font=\sffamily\footnotesize] at (26.0,-6.0) {BC};
  \fill[bkBlue!18] (32.5,-8.6) rectangle (41.5,-3.4); \node[font=\sffamily\footnotesize] at (37.0,-6.0) {BD};
  \fill[bkBlue!18] (43.5,-8.6) rectangle (52.5,-3.4); \node[font=\sffamily\footnotesize] at (48.0,-6.0) {BE};
  \fill[bkOrange!18] (-0.5,-14.6) rectangle (8.5,-9.4); \node[font=\sffamily\footnotesize] at (4.0,-12.0) {CA};
  \fill[bkOrange!18] (10.5,-14.6) rectangle (19.5,-9.4); \node[font=\sffamily\footnotesize] at (15.0,-12.0) {CB};
  \fill[bkBlue!18] (32.5,-14.6) rectangle (41.5,-9.4); \node[font=\sffamily\footnotesize] at (37.0,-12.0) {CD};
  \fill[bkBlue!18] (43.5,-14.6) rectangle (52.5,-9.4); \node[font=\sffamily\footnotesize] at (48.0,-12.0) {CE};
  \fill[bkOrange!18] (-0.5,-20.6) rectangle (8.5,-15.4); \node[font=\sffamily\footnotesize] at (4.0,-18.0) {DA};
  \fill[bkOrange!18] (10.5,-20.6) rectangle (19.5,-15.4); \node[font=\sffamily\footnotesize] at (15.0,-18.0) {DB};
  \fill[bkOrange!18] (21.5,-20.6) rectangle (30.5,-15.4); \node[font=\sffamily\footnotesize] at (26.0,-18.0) {DC};
  \fill[bkBlue!18] (43.5,-20.6) rectangle (52.5,-15.4); \node[font=\sffamily\footnotesize] at (48.0,-18.0) {DE};
  \fill[bkOrange!18] (-0.5,-26.6) rectangle (8.5,-21.4); \node[font=\sffamily\footnotesize] at (4.0,-24.0) {EA};
  \fill[bkOrange!18] (10.5,-26.6) rectangle (19.5,-21.4); \node[font=\sffamily\footnotesize] at (15.0,-24.0) {EB};
  \fill[bkOrange!18] (21.5,-26.6) rectangle (30.5,-21.4); \node[font=\sffamily\footnotesize] at (26.0,-24.0) {EC};
  \fill[bkOrange!18] (32.5,-26.6) rectangle (41.5,-21.4); \node[font=\sffamily\footnotesize] at (37.0,-24.0) {ED};
  \node[font=\sffamily\footnotesize, anchor=west, align=left] at (60,-12) {\R{順番}{じゅん|ばん}をつけて\R{選}{えら}ぶと $5\times4=20$ \R{通}{とお}り。\\\R{青}{あお}の「AB」と、だいだいの「BA」は\\\R{同}{おな}じ\R{組}{くみ}なので、2\R{回}{かい}ずつ\R{数}{かぞ}えている。\\\R{組}{くみ}は $20\div2=10$ \R{通}{とお}り（\R{青}{あお}の10こ）。};
\end{tikzpicture}
\bkfigcaption{5\R{枚}{まい}（A〜E）から2\R{枚}{まい}を\R{選}{えら}ぶ。}\label{bk:3-17:fig-pairs}
\end{bkfigure}
パスカルの\R{三角形}{さん|かく|けい}の5\R{枚}{まい}の\R{段}{だん}（1、5、10、10、5、1）の、\R{左}{ひだり}から2\R{番}{ばん}めの10と\R{合}{あ}っています。3\R{枚}{まい}を\R{選}{えら}ぶなら、\R{選}{えら}ぶ\R{順番}{じゅん|ばん}は $5\times4\times3$ \R{通}{とお}り、3\R{枚}{まい}の\R{並}{なら}べ\R{方}{かた}は $3\times2\times1=6$ \R{通}{とお}りなので、${}_5\mathrm{C}_3=\dfrac{5\times4\times3}{3\times2\times1}=\dfrac{60}{6}=10$ です。

\bkthrrstep{10\R{枚}{まい}\R{投}{な}げたら}
パスカルの\R{三角形}{さん|かく|けい}をあと2\R{段}{だん}のばすと、10\R{枚}{まい}の\R{段}{だん}は 1、10、45、120、210、252、210、120、45、10、1 です（\R{合計}{ごう|けい} $2^{10}=1024$）。\R{表}{おもて}が5\R{枚}{まい}の\R{確率}{かく|りつ}は $\dfrac{252}{1024}\approx0.25$、7\R{枚}{まい}の\R{確率}{かく|りつ}は $\dfrac{120}{1024}\approx0.12$ です。ミナさんの「10\R{回}{かい}で7\R{回}{かい}」は、だいたい8\R{回}{かい}に1\R{回}{かい}は\R{起}{お}こることで、めずらしくありません。

\noindent\begin{minipage}{\linewidth}\begin{motto}[3]
\textbf{\R{数}{かず}の\R{爆発}{ばく|はつ}：トランプの\R{並}{なら}べ\R{方}{かた}}\par
トランプ52\R{枚}{まい}を1\R{列}{れつ}に\R{並}{なら}べる\R{並}{なら}べ\R{方}{かた}は\R{何通}{なん|とお}りでしょう。1\R{枚}{まい}めは52\R{通}{とお}り、2\R{枚}{まい}めは\R{残}{のこ}りの51\R{通}{とお}り、3\R{枚}{まい}めは50\R{通}{とお}り、…と、\R{最後}{さい|ご}の1\R{枚}{まい}まで\R{選}{えら}ぶので、
\bkthrrline{$52\times51\times50\times\cdots\times2\times1$ \R{通}{とお}り}
です。\R{計算}{けい|さん}すると、およそ $8\times10^{67}$ \R{通}{とお}り（8のあとに0が67\R{個}{こ}）になります。\R{地球}{ち|きゅう}\R{上}{うえ}の\R{砂}{すな}つぶの\R{数}{かず}（およそ $10^{19}$ \R{個}{こ}ともいわれます）より、はるかに\R{大}{おお}きな\R{数}{かず}です。よく\R{切}{き}ったトランプの\R{並}{なら}び\R{方}{かた}は、たぶん\R{人類}{じん|るい}の\R{歴史}{れき|し}で\R{一度}{いち|ど}も\R{出}{で}たことのない\R{並}{なら}び\R{方}{かた}なのです。このような「1から $n$ までの\R{数}{かず}を\R{全部}{ぜん|ぶ}かけた\R{数}{かず}」は、5-17で $n!$ という\R{記号}{き|ごう}で\R{書}{か}きます。\R{物理}{ぶつ|り}では、\R{気}{き}\R{体}{からだ}の\R{分子}{ぶん|し}の\R{並}{なら}び\R{方}{かた}の\R{数}{かず}を\R{数}{かぞ}えるときに、この\R{大}{おお}きな\R{数}{かず}が\R{出}{で}てきます。
\end{motto}
\end{minipage}


\bksubsection{4}{たくさんやったときの\R{平均}{へい|きん}：\R{期待値}{き|たい|ち}}

\leavevmode\label{bk:3-17:kitaichi}サイコロを\R{何回}{なん|かい}もふって、\R{出}{で}た\R{目}{め}の\R{平均}{へい|きん}をとると、いくつに\R{近}{ちか}づくでしょう。1から6の\R{目}{め}は、どれも $\dfrac{1}{6}$ ずつ\R{出}{で}ます。600\R{回}{かい}ふれば、どの\R{目}{め}もおよそ100\R{回}{かい}ずつ\R{出}{で}るので、\R{目}{め}の\R{合計}{ごう|けい}はおよそ
\[ 1\times100+2\times100+3\times100+4\times100+5\times100+6\times100=2100 \]
で、\R{平均}{へい|きん}は $2100\div600=3.5$ です。これを1\R{回}{かい}あたりで\R{書}{か}くと、
\[ 1\times\frac{1}{6}+2\times\frac{1}{6}+3\times\frac{1}{6}+4\times\frac{1}{6}+5\times\frac{1}{6}+6\times\frac{1}{6}=\frac{21}{6}=3.5 \]
となります。このように、（\R{値}{あたい}×その\R{確率}{かく|りつ}）を\R{全部}{ぜん|ぶ}たした\R{数}{かず}を\textbf{\R{期待値}{き|たい|ち}}といいます。\bkidx{きたいち}{\R{期待値}{き|たい|ち}}たくさんくり\R{返}{かえ}したときの、1\R{回}{かい}あたりの\R{平均}{へい|きん}です。3.5という\R{目}{め}はありませんが、\R{期待値}{き|たい|ち}は\R{平均}{へい|きん}なので、ない\R{目}{め}でもかまいません。

コイン10\R{枚}{まい}で\R{表}{おもて}が\R{出}{で}る\R{枚数}{まい|すう}の\R{期待値}{き|たい|ち}は、1\R{枚}{まい}あたり $\dfrac{1}{2}$ \R{枚}{まい}ずつなので、$10\times\dfrac{1}{2}=5$ \R{枚}{まい}です。パスカルの\R{三角形}{さん|かく|けい}の10\R{枚}{まい}の\R{段}{だん}で、いちばん\R{大}{おお}きい252がまん\R{中}{なか}の5\R{枚}{まい}にあるのと\R{合}{あ}っています。

\noindent\begin{minipage}{\linewidth}\begin{motto}[3]
\textbf{このコインは、いかさま？（\R{仮説検定}{か|せつ|けん|てい}）}\par
コインを10\R{回}{かい}\R{投}{な}げたら、\R{表}{おもて}が9\R{回}{かい}\R{以上}{い|じょう}\R{出}{で}ました。「このコインは\R{正}{ただ}しい（\R{表}{おもて}の\R{確率}{かく|りつ} $\dfrac{1}{2}$）」と\R{仮}{かり}に\R{考}{かんが}えると、9\R{回}{かい}\R{以上}{い|じょう}\R{表}{おもて}になるのは、パスカルの\R{三角形}{さん|かく|けい}の10\R{枚}{まい}の\R{段}{だん}から $\dfrac{10+1}{1024}=\dfrac{11}{1024}\approx0.011$、\R{約}{やく}1\%です。「100\R{回}{かい}に1\R{回}{かい}しか\R{起}{お}きないことが\R{起}{お}きた」ので、「\R{正}{ただ}しいコインだ」という\R{仮}{かり}の\R{考}{かんが}えのほうを\R{疑}{うたが}います。このように\R{確率}{かく|りつ}を\R{使}{つか}って\R{判断}{はん|だん}する\R{方法}{ほう|ほう}を、\textbf{\R{仮説検定}{か|せつ|けん|てい}}といいます。\R{工場}{こう|じょう}の\R{品質}{ひん|しつ}の\R{判断}{はん|だん}や、\R{薬}{くすり}の\R{効}{き}き\R{目}{め}を\R{確}{たし}かめるときに\R{使}{つか}われます。
\end{motto}
\end{minipage}

\bksubsection{5}{たくさん\R{足}{た}すと、つりがね\R{形}{かたち}：\R{正規分布}{せい|き|ぶん|ぷ}}

\leavevmode\label{bk:3-17:seiki}\bkhatenaref にもどります。サイコロの\R{目}{め}の\R{和}{わ}の、それぞれの\R{確率}{かく|りつ}をグラフにすると、\R{図}{ず}\ref{bk:3-17:fig-dice}のようになります（3\R{個}{こ}まではすべての\R{場合}{ば|あい}を\R{数}{かぞ}えて、10\R{個}{こ}は\R{表計算}{ひょう|けい|さん}で\R{計算}{けい|さん}しました）。

\begin{bkfigure}
\bkthrrbars{1\R{個}{こ}}{(1,1) (2,1) (3,1) (4,1) (5,1) (6,1)}{0.4}{6.6}{1.3}\hfill
\bkthrrbars{2\R{個}{こ}の\R{和}{わ}}{(2,1) (3,2) (4,3) (5,4) (6,5) (7,6) (8,5) (9,4) (10,3) (11,2) (12,1)}{1.4}{12.6}{7}\hfill
\bkthrrbars{3\R{個}{こ}の\R{和}{わ}}{(3,1) (4,3) (5,6) (6,10) (7,15) (8,21) (9,25) (10,27) (11,27) (12,25) (13,21) (14,15) (15,10) (16,6) (17,3) (18,1)}{2.4}{18.6}{31}\hfill
\bkthrrbars{10\R{個}{こ}の\R{和}{わ}}{(20,0.0014) (21,0.0024) (22,0.004) (23,0.0063) (24,0.0095) (25,0.0137) (26,0.019) (27,0.0254) (28,0.0326) (29,0.0405) (30,0.0485) (31,0.0561) (32,0.0629) (33,0.0682) (34,0.0715) (35,0.0727) (36,0.0715) (37,0.0682) (38,0.0629) (39,0.0561) (40,0.0485) (41,0.0405) (42,0.0326) (43,0.0254) (44,0.019) (45,0.0137) (46,0.0095) (47,0.0063) (48,0.004) (49,0.0024) (50,0.0014)}{19}{51}{0.085}
\bkfigcaption{サイコロの\R{目}{め}の\R{和}{わ}の\R{分布}{ぶん|ぷ}。\R{左}{ひだり}から、1\R{個}{こ}（\R{平}{たい}ら）、2\R{個}{こ}（\R{三角}{さん|かく}）、3\R{個}{こ}（\R{丸}{まる}みのある\R{山}{やま}）、10\R{個}{こ}（つりがね\R{形}{かたち}）。}\label{bk:3-17:fig-dice}
\end{bkfigure}

\textbf{なぜ\R{山}{やま}になるのか}\quad \R{和}{わ}がはしの\R{数}{かず}（10\R{個}{こ}なら10や60）になるのは、\R{全部}{ぜん|ぶ}1の\R{目}{め}、\R{全部}{ぜん|ぶ}6の\R{目}{め}のような、たった1\R{通}{とお}りの\R{作}{つく}り\R{方}{かた}しかありません。まん\R{中}{なか}の35のあたりは、\R{大}{おお}きい\R{目}{め}と\R{小}{ちい}さい\R{目}{め}がうまく\R{打}{う}ち\R{消}{け}し\R{合}{あ}う\R{組}{く}み\R{合}{あ}わせが、とてもたくさんあります。パスカルの\R{三角形}{さん|かく|けい}で、まん\R{中}{なか}の\R{数}{かず}ほど\R{大}{おお}きいのと\R{同}{おな}じ\R{理由}{り|ゆう}です。\textbf{\R{作}{つく}り\R{方}{かた}の\R{数}{かず}が\R{多}{おお}い\R{和}{わ}ほど、よく\R{出}{で}る}のです。

\bkhatenaref の\R{答}{こた}え：1\R{個}{こ}1\R{個}{こ}のサイコロは\R{平}{たい}らでも、たくさん\R{足}{た}すと、\R{作}{つく}り\R{方}{かた}の\R{多}{おお}いまん\R{中}{なか}が\R{高}{たか}くなり、\R{山}{やま}の\R{形}{かたち}になります。\R{足}{た}す\R{数}{かず}をふやすほど、\R{山}{やま}はなめらかな\textbf{つりがね\R{形}{かたち}}に\R{近}{ちか}づきます。

\bkthrrstep{つりがね\R{形}{かたち}には\R{名前}{な|まえ}がある}
たくさんの\R{偶然}{ぐう|ぜん}のずれが\R{足}{た}し\R{合}{あ}わさってできるちらばりは、どんなもとのちらばり\R{方}{かた}から\R{始}{はじ}めても、いつも\R{同}{おな}じ\R{形}{かたち}のつりがね\R{形}{かたち}に\R{近}{ちか}づきます。この\R{形}{かたち}のちらばり\R{方}{かた}を\textbf{\R{正規分布}{せい|き|ぶん|ぷ}}といいます。\bkidx{せいきぶんぷ}{\R{正規分布}{せい|き|ぶん|ぷ}}「\R{正規}{せい|き}」は「ふつうの、よく\R{見}{み}かける」という\R{意味}{い|み}です。はかった\R{値}{あたい}のずれ（2-6）も、\R{小}{ちい}さな\R{原因}{げん|いん}がたくさん\R{足}{た}し\R{合}{あ}わさってできるので、\R{正規分布}{せい|き|ぶん|ぷ}に\R{近}{ちか}くなることが\R{多}{おお}いのです。

\begin{bkfigure}
\begin{tikzpicture}
\begin{axis}[bkplot, axis lines=left, width=120mm, height=56mm, xmin=-3.6, xmax=3.6, ymin=0, ymax=0.46, ytick=\empty, xtick={-3,-2,-1,0,1,2,3}, xticklabels={$-3\sigma$, $-2\sigma$, $-\sigma$, \R{平均}{へい|きん}, $+\sigma$, $+2\sigma$, $+3\sigma$}, grid=none]
  \addplot[fill=bkBlue!12, draw=none, domain=-2:2, samples=80] {0.3989*exp(-x^2/2)} \closedcycle;
  \addplot[fill=bkBlue!35, draw=none, domain=-1:1, samples=60] {0.3989*exp(-x^2/2)} \closedcycle;
  \addplot[bkBlue!80!black, line width=1.4pt, domain=-3.5:3.5, samples=120] {0.3989*exp(-x^2/2)};
  \draw[{Stealth[length=3pt]}-{Stealth[length=3pt]}, bkInk, line width=0.6pt] (axis cs:-1,0.15) -- (axis cs:1,0.15);
  \node[font=\sffamily\small\bfseries, fill=bkBlue!35, inner sep=1pt] at (axis cs:0,0.12) {\R{約}{やく}68\%};
  \draw[{Stealth[length=3pt]}-{Stealth[length=3pt]}, bkInk, line width=0.6pt] (axis cs:-2,0.03) -- (axis cs:2,0.03);
  \node[font=\sffamily\small\bfseries, anchor=south] at (axis cs:1.55,0.035) {\R{約}{やく}95\%};
\end{axis}
\end{tikzpicture}
\bkfigcaption{\R{正規分布}{せい|き|ぶん|ぷ}。\R{平均}{へい|きん}±$\sigma$ の\R{間}{あいだ}（こい\R{青}{あお}）に\R{約}{やく}68\%、\R{平均}{へい|きん}±$2\sigma$ の\R{間}{あいだ}（うすい\R{青}{あお}もふくむ）に\R{約}{やく}95\% が\R{入}{はい}る。}\label{bk:3-17:fig-seiki}
\end{bkfigure}

\begin{taisetsu}
\R{正規分布}{せい|き|ぶん|ぷ}では\par
\R{平均}{へい|きん}±$\sigma$ の\R{間}{あいだ}に、\R{全体}{ぜん|たい}の\textbf{\R{約}{やく}68\%}（3\R{分}{ぶん}の2くらい）が\R{入}{はい}る\par
\R{平均}{へい|きん}±$2\sigma$ の\R{間}{あいだ}に、\R{全体}{ぜん|たい}の\textbf{\R{約}{やく}95\%}が\R{入}{はい}る
\end{taisetsu}

3-14では「\R{平均}{へい|きん}±$\sigma$ に、だいたい7\R{割}{わり}」と\R{学}{まな}びました。その7\R{割}{わり}が、\R{正規分布}{せい|き|ぶん|ぷ}の\R{約}{やく}68\%です。

\bkthrrstep{サイコロ10\R{個}{こ}で\R{確}{たし}かめる}
10\R{個}{こ}の\R{和}{わ}の\R{平均}{へい|きん}（\R{期待値}{き|たい|ち}）は、1\R{個}{こ}の\R{期待値}{き|たい|ち}3.5の10\R{倍}{ばい}で35です。\R{表計算}{ひょう|けい|さん}で\R{和}{わ}の\R{分散}{ぶん|さん}を\R{計算}{けい|さん}すると（3-14と\R{同}{おな}じく、\R{偏差}{へん|さ}の2\R{乗}{じょう}に\R{確率}{かく|りつ}をかけてたす）、$\sigma\approx5.4$ になります。
\begin{itemize}[leftmargin=1.4\zw, itemsep=1pt, topsep=2pt, label={\color{bkBlue}●}]
\item \R{平均}{へい|きん}±$\sigma$：$35-5.4=29.6$ から $35+5.4=40.4$ まで。\R{和}{わ}が30〜40になる\R{確率}{かく|りつ}は、\R{数}{かぞ}えると\R{約}{やく}68.7\%。
\item \R{平均}{へい|きん}±$2\sigma$：$35-10.8=24.2$ から $35+10.8=45.8$ まで。\R{和}{わ}が25〜45になる\R{確率}{かく|りつ}は、\R{約}{やく}95.0\%。
\end{itemize}
どちらも、68\%・95\% にとても\R{近}{ちか}くなっています。

\begin{motto}[3]
\textbf{10\R{個}{こ}の\R{和}{わ}の $\sigma\approx5.4$ を、\R{手}{て}で\R{計算}{けい|さん}する}\par
サイコロ1\R{個}{こ}の\R{目}{め}の\R{平均}{へい|きん}は3.5です。\R{偏差}{へん|さ}は $-2.5$、$-1.5$、$-0.5$、0.5、1.5、2.5 で、2\R{乗}{じょう}は6.25、2.25、0.25、0.25、2.25、6.25。どれも\R{同}{おな}じ\R{確率}{かく|りつ}で\R{出}{で}るので、\R{分散}{ぶん|さん}はその\R{平均}{へい|きん}で
\bkthrrline{$(6.25+2.25+0.25+0.25+2.25+6.25)\div6=17.5\div6=2.916\cdots$}
です（3-14と\R{同}{おな}じく、\R{個数}{こ|すう}でわる）。べつべつにふったサイコロの\R{和}{わ}では、\textbf{\R{分散}{ぶん|さん}がたし\R{算}{ざん}になる}ことが\R{知}{し}られています（くわしくは\R{大学}{だい|がく}で）。10\R{個}{こ}なら\R{分散}{ぶん|さん}は10\R{倍}{ばい}の29.16…で、$\sigma=\sqrt{29.16\cdots}\approx5.40$ です。\R{平均}{へい|きん}（\R{期待値}{き|たい|ち}）が10\R{倍}{ばい}になるのに、$\sigma$ は $\sqrt{10}\approx3.16$ \R{倍}{ばい}にしかなりません。だから、たくさん\R{足}{た}すほど、\R{平均}{へい|きん}にくらべてちらばりは\R{小}{ちい}さく\R{見}{み}えます（3-14の $\dfrac{\sigma}{\sqrt{n}}$ と\R{同}{おな}じしくみ）。
\end{motto}

\begin{column}{ガルトンボード：くぎの\R{山}{やま}で\R{正規分布}{せい|き|ぶん|ぷ}を\R{作}{つく}る}
\noindent\begin{minipage}[c]{0.52\linewidth}
\R{板}{いた}にくぎを\R{三角形}{さん|かく|けい}に\R{打}{う}ち、\R{上}{うえ}から\R{玉}{たま}を\R{落}{お}とします。\R{玉}{たま}は、くぎに\R{当}{あ}たるたびに、\R{左}{ひだり}か\R{右}{みぎ}に\R{半分}{はん|ぶん}ずつの\R{確率}{かく|りつ}で\R{分}{わ}かれます。くぎが6\R{段}{だん}なら、\R{玉}{たま}の\R{右}{みぎ}へ\R{行}{い}った\R{回数}{かい|すう}は、コイン6\R{枚}{まい}の\R{表}{おもて}の\R{枚数}{まい|すう}と\R{同}{おな}じです。\R{下}{した}の7つの\R{部屋}{へ|や}に\R{入}{はい}る\R{玉}{たま}の\R{数}{かず}は、パスカルの\R{三角形}{さん|かく|けい}の6\R{段}{だん}め（1、6、15、20、15、6、1）に\R{比例}{ひ|れい}して、つりがね\R{形}{かたち}にたまります。3-15のコラムのゴルトン（ガルトン）が\R{考}{かんが}えた\R{装置}{そう|ち}です。
\end{minipage}\hfill
\begin{minipage}[c]{0.44\linewidth}\centering
\begin{tikzpicture}[x=1mm, y=1mm]
  \foreach \r in {0,...,5} {\foreach \c in {0,...,\r} {{\fill[bkInk] ({25+6*(\c-\r/2)},{40-5.2*\r}) circle (0.7);}}}
  \fill[bkOrange] (25,46) circle (1.3);
  \draw[bkOrange, line width=0.5pt, dash pattern=on 1pt off 1pt] (25,44.5) -- (25,41.2) (25,40) -- (22,35) -- (25,29.6) -- (28,24.4) -- (31,19.2) -- (28,14);
  \foreach \k in {0,...,7} {\draw[bkGray, line width=0.6pt] ({4+6*\k},11) -- ({4+6*\k},0);}
  \draw[bkGray, line width=0.6pt] (4,0) -- (46,0);
  \foreach \k/\h in {0/0.4,1/2.4,2/6,3/8,4/6,5/2.4,6/0.4} {\fill[bkBlue!60] ({4.6+6*\k},0.3) rectangle ({9.4+6*\k},{0.3+1.25*\h});}
  \foreach \k/\v in {0/1,1/6,2/15,3/20,4/15,5/6,6/1} {\node[font=\sffamily\scriptsize, anchor=north] at ({7+6*\k},-0.8) {\v};}
\end{tikzpicture}
\end{minipage}
\end{column}

\bksubsection{6}{\R{全部}{ぜん|ぶ}を\R{調}{しら}べずに、\R{一部}{いち|ぶ}から\R{全体}{ぜん|たい}を\R{知}{し}る：\R{標本調査}{ひょう|ほん|ちょう|さ}}

\leavevmode\label{bk:3-17:hyohon}\R{工場}{こう|じょう}で\R{作}{つく}った\R{電球}{でん|きゅう}が、\R{何時間}{なん|じ|かん}もつかを\R{知}{し}りたいとします。\R{全部}{ぜん|ぶ}の\R{電球}{でん|きゅう}を\R{切}{き}れるまで\R{点}{てん}けて\R{調}{しら}べたら、\R{売}{う}る\R{電球}{でん|きゅう}がなくなってしまいます。そこで、\R{一部}{いち|ぶ}だけを\R{取}{と}り\R{出}{だ}して\R{調}{しら}べ、そこから\R{全体}{ぜん|たい}のようすを\R{予想}{よ|そう}します。

\R{全部}{ぜん|ぶ}を\R{調}{しら}べることを\textbf{\R{全数調査}{ぜん|すう|ちょう|さ}}、\R{一部}{いち|ぶ}を\R{取}{と}り\R{出}{だ}して\R{調}{しら}べることを\textbf{\R{標本調査}{ひょう|ほん|ちょう|さ}}といいます。\bkidx{ひょうほんちょうさ}{\R{標本調査}{ひょう|ほん|ちょう|さ}}\R{調}{しら}べたい\R{全体}{ぜん|たい}を\textbf{\R{母集団}{ぼ|しゅう|だん}}、\R{取}{と}り\R{出}{だ}した\R{一部}{いち|ぶ}を\textbf{\R{標本}{ひょう|ほん}}といいます。\R{国勢調査}{こく|せい|ちょう|さ}（\R{国}{くに}に\R{住}{す}む\R{全員}{ぜん|いん}を\R{調}{しら}べる）は\R{全数調査}{ぜん|すう|ちょう|さ}、テレビの\R{視聴率}{し|ちょう|りつ}やニュースの\R{世論調査}{よ|ろん|ちょう|さ}は\R{標本調査}{ひょう|ほん|ちょう|さ}です。

\bkthrrstep{かたよらないように\R{取}{と}り\R{出}{だ}す}
\R{標本}{ひょう|ほん}は、\R{母集団}{ぼ|しゅう|だん}の\R{中}{なか}から\textbf{かたよりなく}\R{選}{えら}ばないと、\R{全体}{ぜん|たい}を\R{正}{ただ}しく\R{表}{あらわ}しません。たとえば、\R{工場}{こう|じょう}の\R{朝}{あさ}いちばんに\R{作}{つく}った\R{製品}{せい|ひん}だけを\R{調}{しら}べると、\R{機械}{き|かい}が\R{温}{あたた}まる\R{前}{まえ}の\R{不良品}{ふ|りょう|ひん}ばかりかもしれません。くじ\R{引}{ひ}きや\R{乱数}{らん|すう}（\R{表計算}{ひょう|けい|さん}の \texttt{RANDBETWEEN} など）で、どれも\R{同}{おな}じ\R{確率}{かく|りつ}で\R{選}{えら}ばれるようにして\R{取}{と}り\R{出}{だ}します。これを\textbf{\R{無作為}{む|さく|い}に\R{抽出}{ちゅう|しゅつ}する}といいます。

\bkthrrstep{\R{標本}{ひょう|ほん}から\R{全体}{ぜん|たい}を\R{予想}{よ|そう}する}
10000\R{個}{こ}の\R{製品}{せい|ひん}から、\R{無作為}{む|さく|い}に200\R{個}{こ}を\R{取}{と}り\R{出}{だ}して\R{調}{しら}べると、\R{不良品}{ふ|りょう|ひん}が4\R{個}{こ}ありました。\R{標本}{ひょう|ほん}の\R{不良品}{ふ|りょう|ひん}の\R{割合}{わり|あい}は $4\div200=0.02$ です。\R{母集団}{ぼ|しゅう|だん}の\R{割合}{わり|あい}もほぼ\R{同}{おな}じと\R{考}{かんが}えて、\R{全体}{ぜん|たい}の\R{不良品}{ふ|りょう|ひん}はおよそ
\bkthrrline{$10000\times0.02=200$（\R{個}{こ}）}
と\R{予想}{よ|そう}します。

\textbf{\R{標本}{ひょう|ほん}が\R{大}{おお}きいほど、\R{予想}{よ|そう}のずれは\R{小}{ちい}さい}\quad \R{標本}{ひょう|ほん}の\R{割合}{わり|あい}も、\R{偶然}{ぐう|ぜん}でゆれます。コインの\R{割合}{わり|あい}が、\R{回数}{かい|すう}をふやすと0.5に\R{落}{お}ち\R{着}{つ}いたのと\R{同}{おな}じで（①）、\R{標本}{ひょう|ほん}を\R{大}{おお}きくするほど、ゆれは\R{小}{ちい}さくなります。3-14で、$n$ \R{回}{かい}はかって\R{平均}{へい|きん}すると、\R{平均}{へい|きん}のずれが $\dfrac{\sigma}{\sqrt{n}}$ になると\R{学}{まな}びました。\R{標本}{ひょう|ほん}を4\R{倍}{ばい}にすると、ずれは $\dfrac{1}{\sqrt{4}}=\dfrac{1}{2}$ になります。



\begin{taisetsu}
\R{確率}{かく|りつ} $P(A)$ ＝（$A$ の\R{起}{お}こる\R{場合}{ば|あい}の\R{数}{かず}）÷（\R{全部}{ぜん|ぶ}の\R{場合}{ば|あい}の\R{数}{かず}）。0\R{以上}{い|じょう}1\R{以下}{い|か}。\par\vspace{2pt}
コイン $n$ \R{枚}{まい}で\R{表}{おもて}が $k$ \R{枚}{まい}の\R{数}{かず}＝パスカルの\R{三角形}{さん|かく|けい}の $n$ \R{段}{だん}めの $k$ \R{番}{ばん}め＝${}_n\mathrm{C}_k$。\par\vspace{2pt}
\R{期待値}{き|たい|ち}＝（\R{値}{あたい}×その\R{確率}{かく|りつ}）を\R{全部}{ぜん|ぶ}たした\R{数}{かず}＝たくさんやったときの\R{平均}{へい|きん}。\par\vspace{2pt}
\R{正規分布}{せい|き|ぶん|ぷ}：\R{平均}{へい|きん}±$\sigma$ に\R{約}{やく}68\%、\R{平均}{へい|きん}±$2\sigma$ に\R{約}{やく}95\%。
\end{taisetsu}

\begin{kotoba}
\small
\bkword{\R{確率}{かく|りつ} $P(A)$}{ことがら $A$ の\R{起}{お}こりやすさ。たくさんくり\R{返}{かえ}したときの\R{割合}{わり|あい}。（$A$ の\R{場合}{ば|あい}の\R{数}{かず}）÷（\R{全部}{ぜん|ぶ}の\R{場合}{ば|あい}の\R{数}{かず}）。}
\bkword{\R{樹形図}{じゅ|けい|ず}}{\R{場合}{ば|あい}が\R{枝分}{えだ|わ}かれしていくようすを、\R{木}{き}の\R{形}{かたち}にかいた\R{図}{ず}。}
\bkword{パスカルの\R{三角形}{さん|かく|けい}・\R{組}{くみ}\R{合}{あ}せ ${}_n\mathrm{C}_r$}{\R{上}{うえ}の\R{左右}{さ|ゆう}2つの\R{数}{かず}をたしてできる\R{三角形}{さん|かく|けい}。$n$ \R{個}{こ}から $r$ \R{個}{こ}を\R{選}{えら}ぶ\R{数}{かず}。}
\bkword{\R{期待値}{き|たい|ち}}{（\R{値}{あたい}×その\R{確率}{かく|りつ}）を\R{全部}{ぜん|ぶ}たした\R{数}{かず}。たくさんやったときの1\R{回}{かい}あたりの\R{平均}{へい|きん}。}
\bkword{\R{正規分布}{せい|き|ぶん|ぷ}}{\R{偶然}{ぐう|ぜん}のずれがたくさん\R{足}{た}し\R{合}{あ}わさってできる、つりがね\R{形}{かたち}のちらばり。}
\bkword{\R{標本調査}{ひょう|ほん|ちょう|さ}}{\R{母集団}{ぼ|しゅう|だん}から\R{無作為}{む|さく|い}に\R{取}{と}り\R{出}{だ}した\R{標本}{ひょう|ほん}を\R{調}{しら}べて、\R{全体}{ぜん|たい}を\R{予想}{よ|そう}する\R{方法}{ほう|ほう}。}
\end{kotoba}

%% ---------------- 例題 ----------------
\bksection{reidai}{\R{例題}{れい|だい}でたしかめよう}

\begin{reidai}[\R{実験}{じっ|けん}から\R{確率}{かく|りつ}を\R{求}{もと}める]
ペットボトルのキャップを\R{投}{な}げて、\R{上}{うえ}\R{向}{む}き（\R{口}{くち}が\R{上}{うえ}）になる\R{回数}{かい|すう}を\R{調}{しら}べると、2000\R{回}{かい}のうち620\R{回}{かい}でした。（1）\R{上}{うえ}\R{向}{む}きになる\R{確率}{かく|りつ}は、およそいくつですか。（2）このキャップをあと5000\R{回}{かい}\R{投}{な}げると、\R{上}{うえ}\R{向}{む}きはおよそ\R{何回}{なん|かい}ですか。\par
\kangaekata たくさんくり\R{返}{かえ}したときの\R{割合}{わり|あい}が、\R{確率}{かく|りつ}です（\R{統計的}{とう|けい|てき}な\R{確率}{かく|りつ}）。その\R{確率}{かく|りつ}を、\R{次}{つぎ}の\R{回数}{かい|すう}にかけて\R{予想}{よ|そう}します。\par
\shiki（1）\nobreak $620\div2000=0.31$\par
\hspace*{2.6\zw}（2）\nobreak $5000\times0.31=1550$\par
\kotae（1）\nobreak およそ0.31\quad（2）\nobreak およそ1550\R{回}{かい}
\end{reidai}

\begin{reidai}[2\R{個}{こ}のサイコロの\R{確率}{かく|りつ}]
\R{大小}{だい|しょう}2\R{個}{こ}のサイコロをふります。（1）\R{目}{め}の\R{和}{わ}が7になる\R{確率}{かく|りつ}、（2）\R{目}{め}の\R{和}{わ}が10\R{以上}{い|じょう}になる\R{確率}{かく|りつ}を\R{求}{もと}めましょう。\par
\kangaekata \R{全部}{ぜん|ぶ}の\R{場合}{ば|あい}は、6×6の\R{表}{ひょう}で $6\times6=36$ \R{通}{とお}り（②の\R{表}{ひょう}）。それぞれの\R{場合}{ば|あい}の\R{数}{かず}を、\R{表}{ひょう}で\R{数}{かぞ}えます。\par
\shiki（1）\nobreak \R{和}{わ}が7は、\R{表}{ひょう}のななめの6\R{通}{とお}り。$\dfrac{6}{36}=\dfrac{1}{6}$\par
\hspace*{2.6\zw}（2）\nobreak \R{和}{わ}が10は3\R{通}{とお}り、11は2\R{通}{とお}り、12は1\R{通}{とお}りで、$3+2+1=6$ \R{通}{とお}り。$\dfrac{6}{36}=\dfrac{1}{6}$\par
\kotae（1）\nobreak $\dfrac{1}{6}$\quad（2）\nobreak $\dfrac{1}{6}$
\end{reidai}

\begingroup\tcbset{equal height group=bkthrrjA}%
\noindent\begin{minipage}[t]{0.49\linewidth}
\begin{jibun}
\R{画}{が}びょうを1500\R{回}{かい}\R{投}{な}げたら、\R{上}{うえ}\R{向}{む}きが930\R{回}{かい}でした。（1）\R{上}{うえ}\R{向}{む}きになる\R{確率}{かく|りつ}は、およそいくつですか。（2）あと1000\R{回}{かい}\R{投}{な}げると、\R{上}{うえ}\R{向}{む}きはおよそ\R{何回}{なん|かい}ですか。
\bkjans{（1）\nobreak $930\div1500=0.62$ で、およそ0.62\quad（2）\nobreak $1000\times0.62=620$ で、およそ620\R{回}{かい}}
\end{jibun}
\end{minipage}\hfill
\begin{minipage}[t]{0.49\linewidth}
\begin{jibun}
\R{大小}{だい|しょう}2\R{個}{こ}のサイコロをふります。（1）\R{目}{め}の\R{和}{わ}が5になる\R{確率}{かく|りつ}、（2）\R{目}{め}の\R{和}{わ}が4\R{以下}{い|か}になる\R{確率}{かく|りつ}を\R{求}{もと}めましょう。
\bkjans{（1）\nobreak 4\R{通}{とお}りで $\dfrac{4}{36}=\dfrac{1}{9}$\quad（2）\nobreak \R{和}{わ}が2、3、4で $1+2+3=6$ \R{通}{とお}り。$\dfrac{6}{36}=\dfrac{1}{6}$}
\end{jibun}
\end{minipage}\par
\endgroup

\begin{reidai}[\R{樹形図}{じゅ|けい|ず}で\R{数}{かぞ}える]
コインを3\R{枚}{まい}\R{投}{な}げます。（1）\R{表}{おもて}がちょうど2\R{枚}{まい}になる\R{確率}{かく|りつ}、（2）\R{少}{すく}なくとも1\R{枚}{まい}\R{表}{おもて}になる\R{確率}{かく|りつ}を\R{求}{もと}めましょう。\par
\kangaekata \R{図}{ず}\ref{bk:3-17:fig-jukei}の\R{樹形図}{じゅ|けい|ず}で、\R{全部}{ぜん|ぶ}の8\R{通}{とお}りから\R{数}{かぞ}えます。（2）は、「\R{全部}{ぜん|ぶ}\R{裏}{うら}」でない\R{場合}{ば|あい}です。\par
\shiki（1）\nobreak おおう、おうお、うおおの3\R{通}{とお}り。$\dfrac{3}{8}$\par
\hspace*{2.6\zw}（2）\nobreak \R{全部}{ぜん|ぶ}\R{裏}{うら}（ううう）は1\R{通}{とお}りなので、それ\R{以外}{い|がい}は $8-1=7$ \R{通}{とお}り。$\dfrac{7}{8}$\par
\kotae（1）\nobreak $\dfrac{3}{8}$\quad（2）\nobreak $\dfrac{7}{8}$
\end{reidai}

\begin{reidai}[パスカルの\R{三角形}{さん|かく|けい}で\R{数}{かぞ}える]
コインを5\R{枚}{まい}\R{投}{な}げます。（1）\R{表}{おもて}がちょうど2\R{枚}{まい}になる\R{確率}{かく|りつ}を\R{求}{もと}めましょう。その\R{場合}{ば|あい}の\R{数}{かず}を ${}_5\mathrm{C}_2$ のかけ\R{算}{ざん}の\R{式}{しき}でも\R{確}{たし}かめましょう。（2）\R{表}{おもて}が3\R{枚}{まい}\R{以上}{い|じょう}になる\R{確率}{かく|りつ}を\R{求}{もと}めましょう。\par
\kangaekata \R{全部}{ぜん|ぶ}の\R{場合}{ば|あい}は $2^5=32$ \R{通}{とお}り。パスカルの\R{三角形}{さん|かく|けい}の5\R{枚}{まい}の\R{段}{だん}（1、5、10、10、5、1）の、\R{左}{ひだり}から2\R{番}{ばん}め（0\R{番}{ばん}めから\R{数}{かぞ}える）が、\R{表}{おもて}2\R{枚}{まい}の\R{数}{かず}です。（2）の「3\R{枚}{まい}\R{以上}{い|じょう}」は、3\R{枚}{まい}・4\R{枚}{まい}・5\R{枚}{まい}の3つの\R{場合}{ば|あい}を\R{合}{あ}わせます。\par
\shiki（1）\nobreak 5\R{枚}{まい}の\R{段}{だん}で\R{表}{おもて}2\R{枚}{まい}は10\R{通}{とお}り。\R{確率}{かく|りつ}は $\dfrac{10}{32}=\dfrac{5}{16}$\par
\hspace*{2.6\zw}\R{確}{たし}かめ：${}_5\mathrm{C}_2=\dfrac{5\times4}{2\times1}=10$\par
\hspace*{2.6\zw}（2）\nobreak 3\R{枚}{まい}は10\R{通}{とお}り、4\R{枚}{まい}は5\R{通}{とお}り、5\R{枚}{まい}は1\R{通}{とお}り\par
\hspace*{2.6\zw}$10+5+1=16$ \R{通}{とお}り。\R{確率}{かく|りつ}は $\dfrac{16}{32}=\dfrac{1}{2}$\par
\kotae（1）\nobreak $\dfrac{5}{16}$\quad（2）\nobreak $\dfrac{1}{2}$
\end{reidai}

\begingroup\tcbset{equal height group=bkthrrjB}%
\noindent\begin{minipage}[t]{0.49\linewidth}
\begin{jibun}
コインを3\R{枚}{まい}\R{投}{な}げます。（1）\R{表}{おもて}がちょうど1\R{枚}{まい}になる\R{確率}{かく|りつ}、（2）3\R{枚}{まい}とも\R{同}{おな}じ\R{面}{めん}になる\R{確率}{かく|りつ}を\R{求}{もと}めましょう。
\bkjans{（1）\nobreak おうう、うおう、ううおの3\R{通}{とお}りで $\dfrac{3}{8}$\quad（2）\nobreak おおお、うううの2\R{通}{とお}りで $\dfrac{2}{8}=\dfrac{1}{4}$}
\end{jibun}
\end{minipage}\hfill
\begin{minipage}[t]{0.49\linewidth}
\begin{jibun}
コインを6\R{枚}{まい}\R{投}{な}げます。（1）\R{表}{おもて}がちょうど3\R{枚}{まい}になる\R{確率}{かく|りつ}を、パスカルの\R{三角形}{さん|かく|けい}で\R{求}{もと}めましょう。${}_6\mathrm{C}_3$ のかけ\R{算}{ざん}の\R{式}{しき}でも\R{確}{たし}かめましょう。（2）\R{表}{おもて}が4\R{枚}{まい}\R{以上}{い|じょう}になる\R{確率}{かく|りつ}を\R{求}{もと}めましょう。
\bkjans{（1）\nobreak 6\R{枚}{まい}の\R{段}{だん}は1、6、15、20、15、6、1で、\R{表}{おもて}3\R{枚}{まい}は20\R{通}{とお}り。$\frac{20}{64}=\frac{5}{16}$\quad（2）\nobreak $15+6+1=22$ \R{通}{とお}りで $\frac{22}{64}=\frac{11}{32}$}
\end{jibun}
\end{minipage}\par
\endgroup

\begin{reidai}[\R{期待値}{き|たい|ち}を\R{求}{もと}める]
100\R{本}{ぽん}のくじがあります。1\R{等}{とう}（500\R{円}{えん}）が2\R{本}{ほん}、2\R{等}{とう}（100\R{円}{えん}）が10\R{本}{ぽん}、3\R{等}{とう}（10\R{円}{えん}）が30\R{本}{ぽん}で、\R{残}{のこ}りははずれです。くじを1\R{本}{ぽん}\R{引}{ひ}くときにもらえる\R{金額}{きん|がく}の\R{期待値}{き|たい|ち}を\R{求}{もと}めましょう。\par
\kangaekata （\R{金額}{きん|がく}×その\R{確率}{かく|りつ}）を\R{全部}{ぜん|ぶ}たします。はずれは0\R{円}{えん}なので、たしても\R{変}{か}わりません。\par
\shiki $500\times\dfrac{2}{100}+100\times\dfrac{10}{100}+10\times\dfrac{30}{100}$\par
\hspace*{2.6\zw}$=10+10+3=23$\par
\kotae 23\R{円}{えん}（このくじを\R{何回}{なん|かい}も\R{引}{ひ}くと、1\R{回}{かい}あたり\R{平均}{へい|きん}23\R{円}{えん}もらえる）
\end{reidai}

\begin{reidai}[サイコロのゲームの\R{期待値}{き|たい|ち}]
サイコロを1\R{回}{かい}ふって、\R{出}{で}た\R{目}{め}の\R{数}{かず}の10\R{倍}{ばい}の\R{点数}{てん|すう}（1なら10\R{点}{てん}、6なら60\R{点}{てん}）がもらえるゲームがあります。1\R{回}{かい}あたりの\R{点数}{てん|すう}の\R{期待値}{き|たい|ち}を\R{求}{もと}めましょう。\par
\kangaekata 1から6の\R{目}{め}は、どれも\R{確率}{かく|りつ} $\dfrac{1}{6}$ です。（\R{点数}{てん|すう}×その\R{確率}{かく|りつ}）を\R{全部}{ぜん|ぶ}たします。\par
\shiki $10\times\dfrac{1}{6}+20\times\dfrac{1}{6}+30\times\dfrac{1}{6}+40\times\dfrac{1}{6}+50\times\dfrac{1}{6}+60\times\dfrac{1}{6}$\par
\hspace*{2.6\zw}$=\dfrac{10+20+30+40+50+60}{6}=\dfrac{210}{6}=35$\par
\kotae 35\R{点}{てん}（\R{目}{め}の\R{期待値}{き|たい|ち}3.5の10\R{倍}{ばい}になっている）
\end{reidai}

\begingroup\tcbset{equal height group=bkthrrjC}%
\noindent\begin{minipage}[t]{0.49\linewidth}
\begin{jibun}
50\R{本}{ぽん}のくじがあります。1\R{等}{とう}（200\R{円}{えん}）が1\R{本}{ぽん}、2\R{等}{とう}（50\R{円}{えん}）が4\R{本}{ほん}、3\R{等}{とう}（10\R{円}{えん}）が10\R{本}{ぽん}で、\R{残}{のこ}りははずれです。1\R{本}{ぽん}\R{引}{ひ}くときにもらえる\R{金額}{きん|がく}の\R{期待値}{き|たい|ち}を\R{求}{もと}めましょう。
\bkjans{$200\times\dfrac{1}{50}=4$、$50\times\dfrac{4}{50}=4$、$10\times\dfrac{10}{50}=2$。$4+4+2=10$ で10\R{円}{えん}}
\end{jibun}
\end{minipage}\hfill
\begin{minipage}[t]{0.49\linewidth}
\begin{jibun}
サイコロを1\R{回}{かい}ふって、\R{偶数}{ぐう|すう}の\R{目}{め}なら100\R{点}{てん}、\R{奇数}{き|すう}の\R{目}{め}なら0\R{点}{てん}がもらえるゲームがあります。1\R{回}{かい}あたりの\R{点数}{てん|すう}の\R{期待値}{き|たい|ち}を\R{求}{もと}めましょう。
\bkjans{$100\times\dfrac{3}{6}+0\times\dfrac{3}{6}=50$ で50\R{点}{てん}}
\end{jibun}
\end{minipage}\par
\endgroup

\begin{reidai}[\R{正規分布}{せい|き|ぶん|ぷ}の68\%と95\%を\R{使}{つか}う]
ある\R{工場}{こう|じょう}の\R{部品}{ぶ|ひん}の\R{長}{なが}さは、\R{平均}{へい|きん}\qty{50.0}{mm}、\R{標準偏差}{ひょう|じゅん|へん|さ}\qty{0.1}{mm}の\R{正規分布}{せい|き|ぶん|ぷ}になっています。1000\R{個}{こ}\R{作}{つく}ると、（1）49.9〜\qty{50.1}{mm}の\R{部品}{ぶ|ひん}はおよそ\R{何個}{なん|こ}ですか。（2）49.8〜\qty{50.2}{mm}はおよそ\R{何個}{なん|こ}ですか。（3）\qty{50.2}{mm}より\R{長}{なが}い\R{部品}{ぶ|ひん}はおよそ\R{何個}{なん|こ}ですか。\par
\kangaekata 49.9〜50.1は\R{平均}{へい|きん}±$\sigma$、49.8〜50.2は\R{平均}{へい|きん}±$2\sigma$ です。\R{正規分布}{せい|き|ぶん|ぷ}は\R{左右}{さ|ゆう}\R{対称}{たい|しょう}なので、±$2\sigma$ の\R{外}{そと}の5\%は、\R{右}{みぎ}と\R{左}{ひだり}に\R{半分}{はん|ぶん}ずつあります。\par
\begin{center}
\begin{tikzpicture}
\begin{axis}[bkplot, axis lines=left, width=100mm, height=36mm, xmin=49.65, xmax=50.35, ymin=0, ymax=0.46, ytick=\empty, xtick={49.7,49.8,49.9,50.0,50.1,50.2,50.3}, xticklabel style={/pgf/number format/fixed, /pgf/number format/precision=1, /pgf/number format/fixed zerofill}, xlabel={mm}, grid=none]
  \addplot[fill=bkBlue!12, draw=none, domain=49.8:50.2, samples=60] {0.3989*exp(-((x-50)/0.1)^2/2)} \closedcycle;
  \addplot[fill=bkBlue!35, draw=none, domain=49.9:50.1, samples=40] {0.3989*exp(-((x-50)/0.1)^2/2)} \closedcycle;
  \addplot[fill=bkRed!45, draw=none, domain=50.2:50.33, samples=20] {0.3989*exp(-((x-50)/0.1)^2/2)} \closedcycle;
  \addplot[bkBlue!80!black, line width=1.2pt, domain=49.66:50.34, samples=100] {0.3989*exp(-((x-50)/0.1)^2/2)};
  \node[font=\sffamily\footnotesize\bfseries] at (axis cs:50,0.14) {68\%};
  \node[font=\sffamily\footnotesize, text=bkRed!80!black, anchor=south west] at (axis cs:50.22,0.06) {（3）2.5\%};
\end{axis}
\end{tikzpicture}
\end{center}

\shiki（1）\nobreak $1000\times0.68=680$\quad（2）\nobreak $1000\times0.95=950$\par
\hspace*{2.6\zw}（3）\nobreak $1000-950=50$ \R{個}{こ}が\R{外}{そと}にあり、その\R{半分}{はん|ぶん}で $50\div2=25$\par
\kotae（1）\nobreak およそ680\R{個}{こ}\quad（2）\nobreak およそ950\R{個}{こ}\quad（3）\nobreak およそ25\R{個}{こ}
\end{reidai}

\begin{jibun}
クッキー1まいの\R{重}{おも}さは、\R{平均}{へい|きん}\qty{20}{g}、\R{標準偏差}{ひょう|じゅん|へん|さ}\qty{0.5}{g}の\R{正規分布}{せい|き|ぶん|ぷ}です。2000まい\R{作}{つく}ると、（1）19.5〜\qty{20.5}{g}は、（2）19〜\qty{21}{g}は、（3）\qty{21}{g}より\R{重}{おも}いのは、それぞれおよそ\R{何}{なん}まいですか。
\bkjans{（1）\nobreak $2000\times0.68=1360$ で、およそ1360まい\quad（2）\nobreak $2000\times0.95=1900$ で、およそ1900まい\quad（3）\nobreak $(2000-1900)\div2=50$ で、およそ50まい}
\end{jibun}

\begin{reidai}[\R{標本}{ひょう|ほん}から\R{全体}{ぜん|たい}を\R{予想}{よ|そう}する]
ふくろに\R{白}{しろ}と\R{黒}{くろ}の\R{碁石}{ご|いし}が\R{合}{あ}わせて600\R{個}{こ}\R{入}{はい}っています。よくかき\R{混}{ま}ぜて30\R{個}{こ}\R{取}{と}り\R{出}{だ}すと、\R{黒}{くろ}が12\R{個}{こ}でした。ふくろの\R{中}{なか}の\R{黒}{くろ}は、およそ\R{何個}{なん|こ}と\R{考}{かんが}えられますか。\par
\kangaekata よくかき\R{混}{ま}ぜて\R{取}{と}り\R{出}{だ}したので、\R{無作為}{む|さく|い}な\R{標本}{ひょう|ほん}です。\R{標本}{ひょう|ほん}の\R{黒}{くろ}の\R{割合}{わり|あい}は、ふくろ\R{全体}{ぜん|たい}の\R{黒}{くろ}の\R{割合}{わり|あい}とほぼ\R{同}{おな}じと\R{考}{かんが}えます。\par
\begin{center}
\begin{tikzpicture}[x=1mm, y=1mm]
  \filldraw[fill=bkOrange!15, draw=bkOrange!70!black, line width=0.7pt, rounded corners=3mm] (0,0) -- (36,0) -- (40,16) -- (-4,16) -- cycle;
  \draw[bkOrange!70!black, line width=0.7pt] (-4,16) .. controls (5,19) and (27,19) .. (40,16);
  \node[font=\sffamily\footnotesize, align=center] at (18,8) {600\R{個}{こ}\\（\R{黒}{くろ}は？）};
  \draw[bk arrow, bkGray, line width=0.9pt] (43,8) -- (53,8);
  \node[font=\sffamily\footnotesize, anchor=south] at (48,9.5) {30\R{個}{こ}};
  \filldraw[fill=white, draw=bkInk, line width=0.4pt] (58.0,3.0) circle (2);
  \filldraw[fill=white, draw=bkInk, line width=0.4pt] (62.6,3.0) circle (2);
  \filldraw[fill=bkInk, draw=bkInk, line width=0.4pt] (67.2,3.0) circle (2);
  \filldraw[fill=white, draw=bkInk, line width=0.4pt] (71.8,3.0) circle (2);
  \filldraw[fill=bkInk, draw=bkInk, line width=0.4pt] (76.4,3.0) circle (2);
  \filldraw[fill=white, draw=bkInk, line width=0.4pt] (81.0,3.0) circle (2);
  \filldraw[fill=white, draw=bkInk, line width=0.4pt] (85.6,3.0) circle (2);
  \filldraw[fill=white, draw=bkInk, line width=0.4pt] (90.2,3.0) circle (2);
  \filldraw[fill=bkInk, draw=bkInk, line width=0.4pt] (94.8,3.0) circle (2);
  \filldraw[fill=white, draw=bkInk, line width=0.4pt] (99.4,3.0) circle (2);
  \filldraw[fill=white, draw=bkInk, line width=0.4pt] (58.0,8.2) circle (2);
  \filldraw[fill=white, draw=bkInk, line width=0.4pt] (62.6,8.2) circle (2);
  \filldraw[fill=bkInk, draw=bkInk, line width=0.4pt] (67.2,8.2) circle (2);
  \filldraw[fill=white, draw=bkInk, line width=0.4pt] (71.8,8.2) circle (2);
  \filldraw[fill=bkInk, draw=bkInk, line width=0.4pt] (76.4,8.2) circle (2);
  \filldraw[fill=white, draw=bkInk, line width=0.4pt] (81.0,8.2) circle (2);
  \filldraw[fill=bkInk, draw=bkInk, line width=0.4pt] (85.6,8.2) circle (2);
  \filldraw[fill=white, draw=bkInk, line width=0.4pt] (90.2,8.2) circle (2);
  \filldraw[fill=white, draw=bkInk, line width=0.4pt] (94.8,8.2) circle (2);
  \filldraw[fill=white, draw=bkInk, line width=0.4pt] (99.4,8.2) circle (2);
  \filldraw[fill=bkInk, draw=bkInk, line width=0.4pt] (58.0,13.4) circle (2);
  \filldraw[fill=white, draw=bkInk, line width=0.4pt] (62.6,13.4) circle (2);
  \filldraw[fill=bkInk, draw=bkInk, line width=0.4pt] (67.2,13.4) circle (2);
  \filldraw[fill=bkInk, draw=bkInk, line width=0.4pt] (71.8,13.4) circle (2);
  \filldraw[fill=white, draw=bkInk, line width=0.4pt] (76.4,13.4) circle (2);
  \filldraw[fill=white, draw=bkInk, line width=0.4pt] (81.0,13.4) circle (2);
  \filldraw[fill=white, draw=bkInk, line width=0.4pt] (85.6,13.4) circle (2);
  \filldraw[fill=bkInk, draw=bkInk, line width=0.4pt] (90.2,13.4) circle (2);
  \filldraw[fill=bkInk, draw=bkInk, line width=0.4pt] (94.8,13.4) circle (2);
  \filldraw[fill=bkInk, draw=bkInk, line width=0.4pt] (99.4,13.4) circle (2);
  \node[font=\sffamily\footnotesize, anchor=west] at (104,8) {\R{黒}{くろ}12\R{個}{こ}};
\end{tikzpicture}
\end{center}
\shiki \R{標本}{ひょう|ほん}の\R{黒}{くろ}の\R{割合}{わり|あい}：$12\div30=0.4$\par
\hspace*{2.6\zw}\R{全体}{ぜん|たい}の\R{黒}{くろ}：$600\times0.4=240$\par
\kotae およそ240\R{個}{こ}
\end{reidai}

\begin{jibun}
\R{工場}{こう|じょう}で\R{作}{つく}ったねじ8000\R{本}{ぼん}から、\R{無作為}{む|さく|い}に400\R{本}{ぽん}を\R{取}{と}り\R{出}{だ}して\R{調}{しら}べると、\R{不良品}{ふ|りょう|ひん}が6\R{本}{ぽん}ありました。8000\R{本}{ぼん}の\R{中}{なか}の\R{不良品}{ふ|りょう|ひん}は、およそ\R{何本}{なん|ぼん}と\R{考}{かんが}えられますか。
\bkjans{$6\div400=0.015$。$8000\times0.015=120$ で、およそ120\R{本}{ぽん}}
\end{jibun}


\begin{reidai}[\R{印}{しるし}をつけて、\R{全体}{ぜん|たい}の\R{数}{かず}を\R{予想}{よ|そう}する]
\R{池}{いけ}の\R{魚}{さかな}の\R{数}{かず}を\R{調}{しら}べます。まず50ぴきをつかまえて\R{印}{しるし}をつけ、\R{池}{いけ}にもどしました。\R{数日後}{すう|じつ|ご}、よく\R{泳}{およ}ぎ\R{回}{かい}ったあとで40ぴきをつかまえると、そのうち4ひきに\R{印}{しるし}がついていました。\R{池}{いけ}の\R{魚}{さかな}は、およそ\R{何}{なん}びきと\R{考}{かんが}えられますか。\par
\begin{center}
\begin{tikzpicture}[x=1mm, y=1mm]
  \filldraw[fill=bkBlue!8, draw=bkBlue!60, line width=0.6pt, rounded corners=4mm] (0,0) rectangle (34,24);
  \node[font=\sffamily\footnotesize, align=center] at (17,12) {\R{池}{いけ}\\（\R{全部}{ぜん|ぶ}で\\\R{何}{なん}びき？）};
  \draw[bk arrow, bkGray, line width=0.9pt] (36,12) -- (44,12);
  \fill[bkBlue!45] (48.0,4.5) ellipse (1.7 and 1); \fill[bkBlue!45] (49.4,4.5) -- (50.8,5.5) -- (50.8,3.5) -- cycle;
  \fill[bkBlue!45] (53.0,4.5) ellipse (1.7 and 1); \fill[bkBlue!45] (54.4,4.5) -- (55.8,5.5) -- (55.8,3.5) -- cycle;
  \fill[bkBlue!45] (58.0,4.5) ellipse (1.7 and 1); \fill[bkBlue!45] (59.4,4.5) -- (60.8,5.5) -- (60.8,3.5) -- cycle;
  \fill[bkRed] (63.0,4.5) ellipse (1.7 and 1); \fill[bkRed] (64.4,4.5) -- (65.8,5.5) -- (65.8,3.5) -- cycle;
  \fill[bkBlue!45] (68.0,4.5) ellipse (1.7 and 1); \fill[bkBlue!45] (69.4,4.5) -- (70.8,5.5) -- (70.8,3.5) -- cycle;
  \fill[bkBlue!45] (73.0,4.5) ellipse (1.7 and 1); \fill[bkBlue!45] (74.4,4.5) -- (75.8,5.5) -- (75.8,3.5) -- cycle;
  \fill[bkBlue!45] (78.0,4.5) ellipse (1.7 and 1); \fill[bkBlue!45] (79.4,4.5) -- (80.8,5.5) -- (80.8,3.5) -- cycle;
  \fill[bkBlue!45] (83.0,4.5) ellipse (1.7 and 1); \fill[bkBlue!45] (84.4,4.5) -- (85.8,5.5) -- (85.8,3.5) -- cycle;
  \fill[bkBlue!45] (88.0,4.5) ellipse (1.7 and 1); \fill[bkBlue!45] (89.4,4.5) -- (90.8,5.5) -- (90.8,3.5) -- cycle;
  \fill[bkBlue!45] (93.0,4.5) ellipse (1.7 and 1); \fill[bkBlue!45] (94.4,4.5) -- (95.8,5.5) -- (95.8,3.5) -- cycle;
  \fill[bkBlue!45] (48.0,9.5) ellipse (1.7 and 1); \fill[bkBlue!45] (49.4,9.5) -- (50.8,10.5) -- (50.8,8.5) -- cycle;
  \fill[bkBlue!45] (53.0,9.5) ellipse (1.7 and 1); \fill[bkBlue!45] (54.4,9.5) -- (55.8,10.5) -- (55.8,8.5) -- cycle;
  \fill[bkBlue!45] (58.0,9.5) ellipse (1.7 and 1); \fill[bkBlue!45] (59.4,9.5) -- (60.8,10.5) -- (60.8,8.5) -- cycle;
  \fill[bkBlue!45] (63.0,9.5) ellipse (1.7 and 1); \fill[bkBlue!45] (64.4,9.5) -- (65.8,10.5) -- (65.8,8.5) -- cycle;
  \fill[bkBlue!45] (68.0,9.5) ellipse (1.7 and 1); \fill[bkBlue!45] (69.4,9.5) -- (70.8,10.5) -- (70.8,8.5) -- cycle;
  \fill[bkBlue!45] (73.0,9.5) ellipse (1.7 and 1); \fill[bkBlue!45] (74.4,9.5) -- (75.8,10.5) -- (75.8,8.5) -- cycle;
  \fill[bkBlue!45] (78.0,9.5) ellipse (1.7 and 1); \fill[bkBlue!45] (79.4,9.5) -- (80.8,10.5) -- (80.8,8.5) -- cycle;
  \fill[bkBlue!45] (83.0,9.5) ellipse (1.7 and 1); \fill[bkBlue!45] (84.4,9.5) -- (85.8,10.5) -- (85.8,8.5) -- cycle;
  \fill[bkBlue!45] (88.0,9.5) ellipse (1.7 and 1); \fill[bkBlue!45] (89.4,9.5) -- (90.8,10.5) -- (90.8,8.5) -- cycle;
  \fill[bkBlue!45] (93.0,9.5) ellipse (1.7 and 1); \fill[bkBlue!45] (94.4,9.5) -- (95.8,10.5) -- (95.8,8.5) -- cycle;
  \fill[bkBlue!45] (48.0,14.5) ellipse (1.7 and 1); \fill[bkBlue!45] (49.4,14.5) -- (50.8,15.5) -- (50.8,13.5) -- cycle;
  \fill[bkBlue!45] (53.0,14.5) ellipse (1.7 and 1); \fill[bkBlue!45] (54.4,14.5) -- (55.8,15.5) -- (55.8,13.5) -- cycle;
  \fill[bkBlue!45] (58.0,14.5) ellipse (1.7 and 1); \fill[bkBlue!45] (59.4,14.5) -- (60.8,15.5) -- (60.8,13.5) -- cycle;
  \fill[bkRed] (63.0,14.5) ellipse (1.7 and 1); \fill[bkRed] (64.4,14.5) -- (65.8,15.5) -- (65.8,13.5) -- cycle;
  \fill[bkBlue!45] (68.0,14.5) ellipse (1.7 and 1); \fill[bkBlue!45] (69.4,14.5) -- (70.8,15.5) -- (70.8,13.5) -- cycle;
  \fill[bkBlue!45] (73.0,14.5) ellipse (1.7 and 1); \fill[bkBlue!45] (74.4,14.5) -- (75.8,15.5) -- (75.8,13.5) -- cycle;
  \fill[bkBlue!45] (78.0,14.5) ellipse (1.7 and 1); \fill[bkBlue!45] (79.4,14.5) -- (80.8,15.5) -- (80.8,13.5) -- cycle;
  \fill[bkBlue!45] (83.0,14.5) ellipse (1.7 and 1); \fill[bkBlue!45] (84.4,14.5) -- (85.8,15.5) -- (85.8,13.5) -- cycle;
  \fill[bkRed] (88.0,14.5) ellipse (1.7 and 1); \fill[bkRed] (89.4,14.5) -- (90.8,15.5) -- (90.8,13.5) -- cycle;
  \fill[bkBlue!45] (93.0,14.5) ellipse (1.7 and 1); \fill[bkBlue!45] (94.4,14.5) -- (95.8,15.5) -- (95.8,13.5) -- cycle;
  \fill[bkBlue!45] (48.0,19.5) ellipse (1.7 and 1); \fill[bkBlue!45] (49.4,19.5) -- (50.8,20.5) -- (50.8,18.5) -- cycle;
  \fill[bkBlue!45] (53.0,19.5) ellipse (1.7 and 1); \fill[bkBlue!45] (54.4,19.5) -- (55.8,20.5) -- (55.8,18.5) -- cycle;
  \fill[bkBlue!45] (58.0,19.5) ellipse (1.7 and 1); \fill[bkBlue!45] (59.4,19.5) -- (60.8,20.5) -- (60.8,18.5) -- cycle;
  \fill[bkBlue!45] (63.0,19.5) ellipse (1.7 and 1); \fill[bkBlue!45] (64.4,19.5) -- (65.8,20.5) -- (65.8,18.5) -- cycle;
  \fill[bkRed] (68.0,19.5) ellipse (1.7 and 1); \fill[bkRed] (69.4,19.5) -- (70.8,20.5) -- (70.8,18.5) -- cycle;
  \fill[bkBlue!45] (73.0,19.5) ellipse (1.7 and 1); \fill[bkBlue!45] (74.4,19.5) -- (75.8,20.5) -- (75.8,18.5) -- cycle;
  \fill[bkBlue!45] (78.0,19.5) ellipse (1.7 and 1); \fill[bkBlue!45] (79.4,19.5) -- (80.8,20.5) -- (80.8,18.5) -- cycle;
  \fill[bkBlue!45] (83.0,19.5) ellipse (1.7 and 1); \fill[bkBlue!45] (84.4,19.5) -- (85.8,20.5) -- (85.8,18.5) -- cycle;
  \fill[bkBlue!45] (88.0,19.5) ellipse (1.7 and 1); \fill[bkBlue!45] (89.4,19.5) -- (90.8,20.5) -- (90.8,18.5) -- cycle;
  \fill[bkBlue!45] (93.0,19.5) ellipse (1.7 and 1); \fill[bkBlue!45] (94.4,19.5) -- (95.8,20.5) -- (95.8,18.5) -- cycle;
  \node[font=\sffamily\footnotesize, anchor=west, align=left] at (100,12) {40ぴき\R{中}{なか}\\{\color{bkRed}\R{印}{しるし}つき}4ひき};
\end{tikzpicture}
\end{center}
\kangaekata 2\R{回}{かい}めにつかまえた40ぴきは、\R{池}{いけ}\R{全体}{ぜん|たい}からの\R{無作為}{む|さく|い}な\R{標本}{ひょう|ほん}と\R{考}{かんが}えます。\R{標本}{ひょう|ほん}の\R{中}{なか}の「\R{印}{しるし}つき」の\R{割合}{わり|あい}は、\R{池}{いけ}\R{全体}{ぜん|たい}の「\R{印}{しるし}つき」の\R{割合}{わり|あい}とほぼ\R{同}{おな}じです。\R{池}{いけ}\R{全体}{ぜん|たい}の\R{印}{しるし}つきは50ぴきとわかっているので、\R{全体}{ぜん|たい}の\R{数}{かず}を $x$ ひきとして、\R{割合}{わり|あい}が\R{等}{ひと}しい\R{式}{しき}を\R{作}{つく}ります（2-5）。\par
\shiki \R{標本}{ひょう|ほん}の\R{印}{しるし}つきの\R{割合}{わり|あい}：$4\div40=0.1$\par
\hspace*{2.6\zw}\R{池}{いけ}\R{全体}{ぜん|たい}：$50\div x=0.1$ だから、$x=50\div0.1=500$\par
\kotae およそ500ぴき（\R{確}{たし}かめ：500ぴきのうち50ぴきが\R{印}{しるし}つきなら、\R{割合}{わり|あい}は $50\div500=0.1$）
\end{reidai}

\begin{jibun}
ある\R{島}{しま}の\R{野}{の}ウサギの\R{数}{かず}を\R{調}{しら}べます。30ぴきをつかまえて\R{印}{しるし}をつけて\R{放}{はな}し、しばらくしてから60ぴきをつかまえると、そのうち5ひきに\R{印}{しるし}がついていました。\R{島}{しま}の\R{野}{の}ウサギは、およそ\R{何}{なん}びきと\R{考}{かんが}えられますか。
\bkjans{\R{印}{しるし}つきの\R{割合}{わり|あい} $5\div60=\frac{1}{12}$ より $x=30\times12=360$。およそ360ぴき}
\end{jibun}

%% ---------------- まちがえやすいところ ----------------
\bksection{tsumazuki}{まちがえやすいところをチェック}

\begin{tsumazuki}{2\R{個}{こ}のサイコロの\R{和}{わ}の\R{確率}{かく|りつ}}
\bkng{\R{和}{わ}は2から12までの11\R{通}{とお}りなので、\R{和}{わ}が7になる\R{確率}{かく|りつ}は $\dfrac{1}{11}$}
\bkok{\R{全部}{ぜん|ぶ}の\R{場合}{ば|あい}は36\R{通}{とお}りで、\R{和}{わ}が7は6\R{通}{とお}り。\R{確率}{かく|りつ}は $\dfrac{6}{36}=\dfrac{1}{6}$}
2から12までの11\R{通}{とお}りは、\R{同様}{どう|よう}に\R{確}{たし}からしくありません。\R{和}{わ}が2になるのは1\R{通}{とお}り、7になるのは6\R{通}{とお}りで、\R{出}{で}やすさがちがいます。\R{確率}{かく|りつ}は、\R{同様}{どう|よう}に\R{確}{たし}からしい\R{場合}{ば|あい}で\R{数}{かぞ}えます。
\end{tsumazuki}

\begin{tsumazuki}{コイン2\R{枚}{まい}で「1\R{枚}{まい}\R{表}{おもて}・1\R{枚}{まい}\R{裏}{うら}」}
\bkng{「2\R{枚}{まい}\R{表}{おもて}」「1\R{枚}{まい}ずつ」「2\R{枚}{まい}\R{裏}{うら}」の3\R{通}{とお}りなので、$\dfrac{1}{3}$}
\bkok{おお・おう・うお・ううの4\R{通}{とお}りのうち2\R{通}{とお}りで、\mbox{$\dfrac{2}{4}=\dfrac{1}{2}$}}
2\R{枚}{まい}のコインを\R{区別}{く|べつ}して\R{数}{かぞ}えます。「1\R{枚}{まい}ずつ」には、「1\R{枚}{まい}めが\R{表}{おもて}」と「2\R{枚}{まい}めが\R{表}{おもて}」の2つの\R{場合}{ば|あい}があります。
\end{tsumazuki}

\begin{tsumazuki}{\R{前}{まえ}の\R{結果}{けっ|か}が、\R{次}{つぎ}に\R{影響}{えい|きょう}する？}
\bkng{5\R{回}{かい}\R{続}{つづ}けて\R{裏}{うら}が\R{出}{で}たから、\R{次}{つぎ}は\R{表}{おもて}が\R{出}{で}やすい}
\bkok{コインには\R{記憶}{き|おく}がない。\R{次}{つぎ}も\R{表}{おもて}の\R{確率}{かく|りつ}は $\dfrac{1}{2}$}
\R{割合}{わり|あい}が0.5に\R{近}{ちか}づくのは、\R{何百回}{なん|びゃっ|かい}もくり\R{返}{かえ}して、はじめのかたよりがうすまるからです。
\end{tsumazuki}

\begin{tsumazuki}{\R{確率}{かく|りつ} $\dfrac{1}{6}$ なら、6\R{回}{かい}で\R{必}{かなら}ず1\R{回}{かい}？}
\bkng{6\R{回}{かい}ふれば、1の\R{目}{め}は\R{必}{かなら}ず1\R{回}{かい}\R{出}{で}る}
\bkok{6\R{回}{かい}ふって、1の\R{目}{め}が1\R{回}{かい}も\R{出}{で}ないこともある}
\R{確率}{かく|りつ}は、たくさんくり\R{返}{かえ}したときの\R{割合}{わり|あい}です。6\R{回}{かい}では、0\R{回}{かい}のことも2\R{回}{かい}\R{以上}{い|じょう}のこともあります。
\end{tsumazuki}

\begin{tsumazuki}{\R{平均}{へい|きん}＋$2\sigma$ より\R{大}{おお}きいのは\R{何}{なん}\%？}
\bkng{±$2\sigma$ の\R{外}{そと}は5\%なので、5\%}
\bkok{\R{外}{そと}の5\%は\R{左右}{さ|ゆう}に\R{半分}{はん|ぶん}ずつなので、$5\div2=2.5$ で2.5\%}
\R{正規分布}{せい|き|ぶん|ぷ}は\R{左右}{さ|ゆう}\R{対称}{たい|しょう}です。\R{外}{そと}の5\%には、\R{小}{ちい}さい\R{側}{がわ}と\R{大}{おお}きい\R{側}{がわ}の\R{両方}{りょう|ほう}がふくまれています（\bkreidairef{7}の（3））。
\end{tsumazuki}

%% ---------------- 物理・工学 ----------------
\begin{butsuri}
\par\addvspace{3pt}\textbf{\R{電気}{でん|き}の\R{雑音}{ざつ|おん}：\R{抵抗}{てい|こう}の\R{中}{なか}の\R{電子}{でん|し}のゆれ}\par
\noindent\begin{minipage}[c]{0.56\linewidth}
\R{電気}{でん|き}を\R{流}{なが}していない\R{抵抗}{てい|こう}でも、\R{中}{なか}の\R{電子}{でん|し}は\R{熱}{ねつ}でたえずでたらめに\R{動}{うご}いています。たくさんの\R{電子}{でん|し}の\R{小}{ちい}さな\R{動}{うご}きが\R{足}{た}し\R{合}{あ}わさって、ごく\R{小}{ちい}さな\R{電圧}{でん|あつ}のゆれ（\R{雑音}{ざつ|おん}）が\R{生}{う}まれます。たくさんの\R{偶然}{ぐう|ぜん}を\R{足}{た}したものなので、\R{雑音}{ざつ|おん}の\R{電圧}{でん|あつ}は\R{平均}{へい|きん}0の\R{正規分布}{せい|き|ぶん|ぷ}になります。センサーで\R{小}{ちい}さな\R{信号}{しん|ごう}をはかるときは、この\R{雑音}{ざつ|おん}の $\sigma$ より\R{信号}{しん|ごう}が\R{十分}{じゅう|ぶん}\R{大}{おお}きくなるように、\R{回路}{かい|ろ}を\R{設計}{せっ|けい}します。
\end{minipage}\hfill
\begin{minipage}[c]{0.4\linewidth}\centering
\begin{tikzpicture}
\begin{axis}[bkplot, axis lines=left, width=52mm, height=34mm, xmin=0, xmax=60, ymin=-3, ymax=3, xtick=\empty, ytick={0}, xlabel={\R{時間}{じ|かん}}, ylabel={\R{電圧}{でん|あつ}}, grid=none]
  \addplot[bkBlue, line width=0.7pt] coordinates {(0,0.3) (1,-0.8) (2,1.1) (3,0.2) (4,-1.5) (5,0.6) (6,1.9) (7,-0.4) (8,-1.1) (9,0.9) (10,0.1) (11,-2.1) (12,0.7) (13,1.3) (14,-0.6) (15,0.4) (16,-0.9) (17,1.6) (18,-0.2) (19,-1.3) (20,0.8) (21,0.2) (22,-0.7) (23,2.2) (24,-0.5) (25,-1.7) (26,0.6) (27,1.0) (28,-0.3) (29,0.1) (30,-1.0) (31,1.4) (32,-0.6) (33,0.3) (34,-2.4) (35,0.9) (36,0.5) (37,-0.8) (38,1.2) (39,-0.1) (40,0.7) (41,-1.4) (42,0.2) (43,1.8) (44,-0.9) (45,-0.3) (46,1.1) (47,-1.6) (48,0.4) (49,0.8) (50,-0.5) (51,1.5) (52,-1.2) (53,0.0) (54,0.6) (55,-0.7) (56,2.0) (57,-1.0) (58,0.3) (59,-0.4) (60,0.9)};
\end{axis}
\end{tikzpicture}
\end{minipage}

\par\addvspace{8pt}\textbf{\R{放射線}{ほう|しゃ|せん}を\R{数}{かぞ}える：めったに\R{起}{お}きないことを\R{数}{かぞ}える}\par
\noindent \R{放射線}{ほう|しゃ|せん}をはかる\R{計数器}{けい|すう|き}は、\R{放射線}{ほう|しゃ|せん}が1つ\R{通}{とお}るたびに「カチッ」と1\R{回}{かい}\R{数}{かぞ}えます。\R{放射線}{ほう|しゃ|せん}がいつ\R{来}{く}るかは\R{偶然}{ぐう|ぜん}で\R{決}{き}まるので、\R{同}{おな}じ\R{場所}{ば|しょ}で1\R{分}{ぶん}ずつはかっても、\R{数}{かず}は\R{毎回}{まい|かい}ちがいます。\R{平均}{へい|きん}して100\R{回数}{かい|すう}えられる\R{場所}{ば|しょ}では、ばらつきの\R{標準偏差}{ひょう|じゅん|へん|さ}はおよそ $\sqrt{100}=10$ \R{回}{かい}になることが\R{知}{し}られています（\R{物理}{ぶつ|り}のきまり）。だから「1\R{分}{ぶん}に100\R{回}{かい}」は「100±10\R{回}{かい}くらい」の\R{意味}{い|み}です。1000\R{回数}{かい|すう}えれば、ばらつきは $\sqrt{1000}\approx32$ \R{回}{かい}で、\R{割合}{わり|あい}にすると\R{約}{やく}3\%まで\R{小}{ちい}さくなります。\R{長}{なが}い\R{時間}{じ|かん}はかるほど、\R{正確}{せい|かく}になるのです。

\end{butsuri}

\begin{butsuri}
\par\addvspace{3pt}\textbf{\R{気}{き}\R{体}{からだ}の\R{分子}{ぶん|し}の\R{速}{はや}さ}\par
\noindent \R{部屋}{へ|や}の\R{空気}{くう|き}の\R{分子}{ぶん|し}は、ぶつかり\R{合}{あ}いながら、いろいろな\R{速}{はや}さで\R{飛}{と}び\R{回}{まわ}っています。1つ1つの\R{分子}{ぶん|し}の\R{速}{はや}さは\R{偶然}{ぐう|ぜん}で\R{決}{き}まりますが、ものすごい\R{数}{かず}（1\,cm$^3$に\R{約}{やく} $2.5\times10^{19}$ \R{個}{こ}）の\R{分子}{ぶん|し}があるので、\R{速}{はや}さのちらばり\R{方}{かた}は、\R{温度}{おん|ど}だけで\R{決}{き}まる、いつも\R{同}{おな}じ\R{形}{かたち}になります。\R{速}{はや}さの\R{平均}{へい|きん}は、\R{室温}{しつ|おん}で\R{秒速}{びょう|そく}500\,mくらいです。\R{物理}{ぶつ|り}シリーズでは、\R{確率}{かく|りつ}を\R{使}{つか}って、このちらばり\R{方}{かた}から\R{温度}{おん|ど}や\R{圧力}{あつ|りょく}を\R{説明}{せつ|めい}します。

\par\addvspace{8pt}\textbf{\R{何回}{なん|かい}も\R{重}{かさ}ねて、\R{雑音}{ざつ|おん}を\R{消}{け}す}\par
\noindent\begin{minipage}[c]{0.5\linewidth}
\R{小}{ちい}さな\R{信号}{しん|ごう}が\R{雑音}{ざつ|おん}にうもれて\R{見}{み}えないとき、\R{同}{おな}じ\R{信号}{しん|ごう}を\R{何回}{なん|かい}もはかって、\R{時刻}{じ|こく}ごとに\R{平均}{へい|きん}します。\R{信号}{しん|ごう}は\R{毎回}{まい|かい}\R{同}{おな}じなので、\R{平均}{へい|きん}しても\R{残}{のこ}ります。\R{雑音}{ざつ|おん}は\R{毎回}{まい|かい}でたらめなので、\R{平均}{へい|きん}すると\R{打}{う}ち\R{消}{け}し\R{合}{あ}って\R{小}{ちい}さくなります。3-14の $\dfrac{\sigma}{\sqrt{n}}$ のとおり、16\R{回}{かい}の\R{平均}{へい|きん}で\R{雑音}{ざつ|おん}は $\dfrac{1}{\sqrt{16}}=\dfrac{1}{4}$ になります。\R{心電図}{しん|でん|ず}や、\R{遠}{とお}くの\R{星}{ほし}からの\R{電波}{でん|ぱ}をとらえる\R{望遠鏡}{ぼう|えん|きょう}などで\R{使}{つか}われています。
\end{minipage}\hfill
\begin{minipage}[c]{0.47\linewidth}\centering
\begin{tikzpicture}
\begin{axis}[bkplot, axis lines=left, name=s1, width=60mm, height=28mm, xmin=0, xmax=80, ymin=-4, ymax=4, xtick=\empty, ytick=\empty, grid=none, title={\sffamily\footnotesize 1\R{回}{かい}だけ}, title style={yshift=-1.5mm}]
  \addplot[bkBlue, line width=0.6pt] coordinates {(0,0.11) (1,1.66) (2,-0.81) (3,1.64) (4,0.28) (5,0.39) (6,3.09) (7,1.08) (8,0.90) (9,1.86) (10,2.35) (11,0.95) (12,1.66) (13,-0.28) (14,0.37) (15,0.18) (16,-1.01) (17,-1.36) (18,-1.64) (19,-0.13) (20,-0.21) (21,-0.54) (22,-0.23) (23,-2.06) (24,-0.68) (25,-0.42) (26,0.09) (27,-1.91) (28,-1.43) (29,-3.41) (30,-1.60) (31,-3.62) (32,-2.65) (33,0.43) (34,-3.45) (35,0.25) (36,-0.19) (37,-0.83) (38,0.24) (39,0.48) (40,1.25) (41,-0.12) (42,-0.40) (43,-0.27) (44,-0.60) (45,0.65) (46,-0.13) (47,2.17) (48,-1.29) (49,-0.32) (50,-0.14) (51,-1.52) (52,3.23) (53,-2.00) (54,0.47) (55,0.08) (56,2.57) (57,-1.93) (58,1.60) (59,-0.72) (60,-0.19) (61,-0.96) (62,0.46) (63,-1.82) (64,-0.68) (65,-0.28) (66,1.40) (67,-3.78) (68,0.88) (69,0.15) (70,-1.58) (71,-0.62) (72,-1.51) (73,1.09) (74,-0.56) (75,-0.97) (76,-0.86) (77,-0.69) (78,-0.52) (79,-1.21) (80,2.47)};
\end{axis}
\begin{axis}[bkplot, axis lines=left, at={(s1.south)}, anchor=north, yshift=-6mm, width=60mm, height=28mm, xmin=0, xmax=80, ymin=-4, ymax=4, xtick=\empty, ytick=\empty, grid=none, title={\sffamily\footnotesize 16\R{回}{かい}の\R{平均}{へい|きん}}, title style={yshift=-1.5mm}]
  \addplot[bkRed, line width=0.8pt] coordinates {(0,-0.05) (1,-0.10) (2,-0.07) (3,0.13) (4,0.76) (5,0.64) (6,0.88) (7,0.78) (8,0.84) (9,0.96) (10,1.36) (11,0.96) (12,0.82) (13,0.87) (14,1.07) (15,0.33) (16,0.45) (17,0.52) (18,0.63) (19,-0.31) (20,0.30) (21,-0.05) (22,0.07) (23,-0.74) (24,-0.80) (25,-0.74) (26,-1.09) (27,-1.51) (28,-1.21) (29,-1.06) (30,-1.13) (31,-1.03) (32,-0.86) (33,-0.76) (34,-1.12) (35,-0.45) (36,-0.51) (37,-0.84) (38,-0.26) (39,0.31) (40,0.51) (41,-0.30) (42,0.58) (43,0.28) (44,0.85) (45,0.63) (46,1.18) (47,1.22) (48,0.66) (49,0.16) (50,0.79) (51,0.73) (52,1.46) (53,1.28) (54,1.32) (55,0.97) (56,0.31) (57,0.86) (58,0.21) (59,0.16) (60,0.10) (61,-0.54) (62,-0.61) (63,-0.65) (64,-0.73) (65,-0.71) (66,-0.75) (67,-0.44) (68,-0.73) (69,-1.43) (70,-1.13) (71,-1.25) (72,-0.89) (73,-0.52) (74,-0.81) (75,-0.66) (76,-0.71) (77,-0.28) (78,-0.74) (79,-0.13) (80,-0.13)};
\end{axis}
\end{tikzpicture}
\end{minipage}
\end{butsuri}

%% ---------------- 章末問題 ----------------
\bkendproblems
\bkproblemgroup{\R{基本}{き|ほん}}{2}

\bkproblem あるボタンを\R{投}{な}げると、1200\R{回}{かい}のうち\R{表}{おもて}\R{向}{む}きが780\R{回}{かい}でした。（1）\R{表}{おもて}\R{向}{む}きになる\R{確率}{かく|りつ}は、およそいくつですか。（2）あと2000\R{回}{かい}\R{投}{な}げると、\R{表}{おもて}\R{向}{む}きはおよそ\R{何回}{なん|かい}ですか。
\bkpans{（1）\nobreak $780\div1200=0.65$ で、およそ0.65\quad（2）\nobreak $2000\times0.65=1300$ で、およそ1300\R{回}{かい}}

\bkproblem \R{大小}{だい|しょう}2\R{個}{こ}のサイコロをふります。（1）\R{目}{め}の\R{和}{わ}が9になる\R{確率}{かく|りつ}、（2）\R{目}{め}の\R{和}{わ}が3\R{以下}{い|か}になる\R{確率}{かく|りつ}を\R{求}{もと}めましょう。
\bkpans{（1）\nobreak 4\R{通}{とお}りで $\dfrac{4}{36}=\dfrac{1}{9}$\quad（2）\nobreak \R{和}{わ}が2、3で $1+2=3$ \R{通}{とお}り。$\dfrac{3}{36}=\dfrac{1}{12}$}

\bkproblem コインを3\R{枚}{まい}\R{投}{な}げます。（1）3\R{枚}{まい}とも\R{表}{おもて}になる\R{確率}{かく|りつ}、（2）\R{少}{すく}なくとも1\R{枚}{まい}\R{裏}{うら}になる\R{確率}{かく|りつ}を\R{求}{もと}めましょう。
\bkpans{（1）\nobreak おおおの1\R{通}{とお}りで $\dfrac{1}{8}$\quad（2）\nobreak \R{全部}{ぜん|ぶ}\R{表}{おもて}\R{以外}{い|がい}の $8-1=7$ \R{通}{とお}りで $\dfrac{7}{8}$}

\bkproblem コインを4\R{枚}{まい}\R{投}{な}げます。\R{表}{おもて}がちょうど2\R{枚}{まい}になる\R{確率}{かく|りつ}を、パスカルの\R{三角形}{さん|かく|けい}で\R{求}{もと}めましょう。${}_4\mathrm{C}_2$ のかけ\R{算}{ざん}の\R{式}{しき}でも\R{確}{たし}かめましょう。
\bkpans{4\R{枚}{まい}の\R{段}{だん}は1、4、6、4、1で、\R{表}{おもて}2\R{枚}{まい}は6\R{通}{とお}り。$\dfrac{6}{16}=\dfrac{3}{8}$。${}_4\mathrm{C}_2=\dfrac{4\times3}{2\times1}=6$}

\bkproblem \R{大小}{だい|しょう}2\R{個}{こ}のサイコロをふります。\R{目}{め}の\R{和}{わ}が11\R{以上}{い|じょう}になる\R{確率}{かく|りつ}を\R{求}{もと}めましょう。
\bkpans{\R{和}{わ}が11は2\R{通}{とお}り、12は1\R{通}{とお}りで、$2+1=3$ \R{通}{とお}り。$\dfrac{3}{36}=\dfrac{1}{12}$}

\bkproblem 200\R{本}{ぽん}のくじがあります。1\R{等}{とう}（1000\R{円}{えん}）が2\R{本}{ほん}、2\R{等}{とう}（100\R{円}{えん}）が20\R{本}{ぽん}、3\R{等}{とう}（20\R{円}{えん}）が50\R{本}{ぽん}で、\R{残}{のこ}りははずれです。1\R{本}{ぽん}\R{引}{ひ}くときにもらえる\R{金額}{きん|がく}の\R{期待値}{き|たい|ち}を\R{求}{もと}めましょう。
\bkpans{$1000\times\dfrac{2}{200}=10$、$100\times\dfrac{20}{200}=10$、$20\times\dfrac{50}{200}=5$。$10+10+5=25$ で25\R{円}{えん}}

\bkproblem ある\R{中学校}{ちゅう|がっ|こう}の1\R{年生}{ねん|せい}500\R{人}{にん}の\R{身長}{しん|ちょう}は、\R{平均}{へい|きん}\qty{152}{cm}、\R{標準偏差}{ひょう|じゅん|へん|さ}\qty{7}{cm}の\R{正規分布}{せい|き|ぶん|ぷ}に\R{近}{ちか}いとします。（1）145〜\qty{159}{cm}の\R{人}{ひと}は、（2）138〜\qty{166}{cm}の\R{人}{ひと}は、それぞれおよそ\R{何人}{なん|にん}ですか。
\bkpans{（1）\nobreak \R{平均}{へい|きん}±$\sigma$ なので $500\times0.68=340$ で、およそ340\R{人}{にん}\quad（2）\nobreak \R{平均}{へい|きん}±$2\sigma$ なので $500\times0.95=475$ で、およそ475\R{人}{にん}}

\bkproblem \R{池}{いけ}の\R{水}{みず}を、よくかき\R{混}{ま}ぜてから\qty{1}{L}くみ\R{取}{と}ると、ミジンコが36ぴきいました。\R{池}{いけ}の\R{水}{みず}は\R{全部}{ぜん|ぶ}で\qty{5000}{L}です。\R{池}{いけ}のミジンコは、およそ\R{何}{なん}びきと\R{考}{かんが}えられますか。
\bkpans{\qty{1}{L}あたり36ぴき。$5000\times36=180000$ で、およそ18\R{万}{まん}びき}

\bkproblemgroup{チャレンジ}{3}

\bkproblem コインを10\R{回}{かい}\R{投}{な}げて、\R{表}{おもて}が8\R{回}{かい}\R{以上}{い|じょう}\R{出}{で}る\R{確率}{かく|りつ}を、パスカルの\R{三角形}{さん|かく|けい}の10\R{枚}{まい}の\R{段}{だん}（1、10、45、120、210、252、210、120、45、10、1）を\R{使}{つか}って\R{求}{もと}めましょう。
\bkpans{8\R{回}{かい}は45\R{通}{とお}り、9\R{回}{かい}は10\R{通}{とお}り、10\R{回}{かい}は1\R{通}{とお}りで、$45+10+1=56$ \R{通}{とお}り。$\dfrac{56}{1024}=\dfrac{7}{128}\approx0.055$（\R{約}{やく}5.5\%）}

\bkproblem サイコロを10\R{個}{こ}ふって\R{目}{め}を\R{足}{た}すことを1000\R{回}{かい}くり\R{返}{かえ}します。\R{和}{わ}は\R{平均}{へい|きん}35、\R{標準偏差}{ひょう|じゅん|へん|さ}\R{約}{やく}5.4の\R{正規分布}{せい|き|ぶん|ぷ}に\R{近}{ちか}くなります。（1）\R{和}{わ}が29.6から40.4のあいだになるのは、およそ\R{何回}{なん|かい}ですか。（2）45.8より\R{大}{おお}きくなるのは、およそ\R{何回}{なん|かい}ですか。
\bkpans{（1）\nobreak \R{平均}{へい|きん}±$\sigma$ なので $1000\times0.68=680$ で、およそ680\R{回}{かい}\quad（2）\nobreak \R{平均}{へい|きん}＋$2\sigma=35+10.8=45.8$。\R{外側}{そと|がわ}の5\%の\R{半分}{はん|ぶん}で $1000\times0.05\div2=25$ で、およそ25\R{回}{かい}}

\begin{furikaeri}
\begin{checklist}
\item \R{表}{ひょう}や\R{樹形図}{じゅ|けい|ず}で\R{数}{かぞ}えて、\R{確率}{かく|りつ}と\R{期待値}{き|たい|ち}を\R{求}{もと}められる。\bkfurref{bk:3-17:kazoeru}
\item パスカルの\R{三角形}{さん|かく|けい}で、\R{表}{おもて}が $k$ \R{枚}{まい}になる\R{数}{かず}を\R{求}{もと}められる。\bkfurref{bk:3-17:pascal}
\item \R{正規分布}{せい|き|ぶん|ぷ}の68\%と95\%を\R{使}{つか}える。\bkfurref{bk:3-17:seiki}
\end{checklist}
\end{furikaeri}

\bkfukidashi{delta}{はかる・\R{表}{ひょう}にする・\R{点}{てん}にする・\R{線}{せん}を\R{引}{ひ}く・ばらつきを\R{考}{かんが}える。\R{第}{だい}3\R{巻}{かん}で、データからきまりを\R{見}{み}つける\R{道具}{どう|ぐ}がそろったね。}
